% ============================================================
% report 学术成果展示 — QUDA multigrid 算法实现（可复现实现手册版）
% 编译: xelatex 两遍；交付前 Overfull=0, Float too large=0
% ============================================================
\documentclass[UTF8,aspectratio=169,11pt]{ctexbeamer}

\usepackage{amsmath,amssymb}
\usepackage{graphicx}
\usepackage{booktabs}
\usepackage{tabularx}
\usepackage{array}
\usepackage{tikz}
\usetikzlibrary{arrows.meta,positioning,calc,fit,shapes.geometric}
\usepackage{listings}
\usepackage{xcolor}
\usepackage{microtype}
\usepackage{hyperref}

% ===== 色板 =====
\definecolor{primary}{HTML}{17365D}
\definecolor{accent}{HTML}{1F77B4}
\definecolor{evidence}{HTML}{1E8449}
\definecolor{warning}{HTML}{C55A11}
\definecolor{muted}{HTML}{5B6573}
\definecolor{panel}{HTML}{F4F7FA}
\definecolor{linegray}{HTML}{CBD5E1}
\definecolor{good}{HTML}{EAF5EE}
\definecolor{warnbg}{HTML}{FFF4E8}
\definecolor{violet}{HTML}{6C4A9E}

\setbeamersize{text margin left=0.55cm,text margin right=0.55cm}
\setlength{\emergencystretch}{2em}
\setbeamertemplate{navigation symbols}{}
\setbeamertemplate{blocks}[rounded][shadow=false]
\setbeamercolor{normal text}{fg=black!86,bg=white}
\setbeamercolor{structure}{fg=primary}
\setbeamercolor{frametitle}{fg=primary,bg=white}
\setbeamercolor{block title}{fg=primary,bg=panel}
\setbeamercolor{block body}{fg=black!86,bg=panel}
\setbeamercolor{alerted text}{fg=warning}
\setbeamerfont{frametitle}{size=\large,series=\bfseries}
\setbeamerfont{normal text}{size=\normalsize}
\setbeamertemplate{frametitle}{\vskip0.12cm\insertframetitle\par\vskip0.08cm}
\setbeamertemplate{footline}{%
  \hbox{\begin{beamercolorbox}[wd=\paperwidth,ht=2.3ex,dp=0.9ex,
    leftskip=0.55cm,rightskip=0.55cm]{author in head/foot}%
    \scriptsize\color{muted}\insertshortauthor\hfill
    \insertshorttitle\hfill\insertframenumber/\inserttotalframenumber
  \end{beamercolorbox}}}

% ===== 来源与证据状态 =====
\newcommand{\source}[1]{\vspace{0.05em}\par{\raggedright\scriptsize\color{muted}来源：#1\par}}
\newcommand{\verified}{\textcolor{evidence}{确证}}
\newcommand{\inferred}{\textcolor{accent}{推断}}
\newcommand{\unverified}{\textcolor{warning}{未验证}}
\newcommand{\codepath}[1]{\nolinkurl{#1}}

\tikzset{
  flow/.style={draw=accent,fill=accent!8,rounded corners=2pt,align=center,
    minimum height=0.72cm,text width=2.05cm,font=\scriptsize},
  flowwide/.style={draw=accent,fill=accent!8,rounded corners=2pt,align=center,
    minimum height=0.68cm,text width=2.55cm,font=\scriptsize},
  state/.style={draw=primary,fill=primary!7,rounded corners=2pt,align=center,
    minimum height=0.68cm,text width=2.55cm,font=\small},
  phase/.style={draw=violet,fill=violet!8,rounded corners=2pt,align=center,
    minimum height=0.62cm,text width=2.30cm,font=\scriptsize},
  decision/.style={draw=warning,fill=warning!10,diamond,aspect=2,align=center,
    inner sep=1.5pt,font=\small},
  arrow/.style={-{Latex[length=2mm]},draw=primary,semithick},
  relation/.style={-{Latex[length=1.7mm]},draw=muted,thin,dashed},
  dotarrow/.style={-{Latex[length=1.7mm]},draw=accent,thin,dotted}
}

\lstset{
  basicstyle=\ttfamily\scriptsize,
  backgroundcolor=\color{panel},frame=single,
  breaklines=true,breakatwhitespace=false,keepspaces=true,
  columns=flexible,firstnumber=auto,
  keywordstyle=\color{primary}\bfseries,
  commentstyle=\color{evidence!70!black},
  stringstyle=\color{accent},
  numbers=left,numberstyle=\tiny\color{muted}
}

\title[QUDA multigrid 复现手册]{QUDA multigrid 算法实现}
\subtitle{从近零模、Galerkin 粗化到 MPI 递归 V-cycle 的可复现逻辑}
\author[源码分析汇报]{源码分析汇报}
\institute{QUDA 1.1.0 参考源码快照}
\date{2026 年 8 月 28 日}

\begin{document}

\begin{frame}[plain]
  \titlepage
\end{frame}

\begin{frame}[t]{本版报告的验收目标：读者可脱离 QUDA 重建同一算法}
  \small
  \begin{columns}[T,totalwidth=\textwidth]
    \begin{column}{0.31\textwidth}
      \begin{block}{数学等价}
        给定同一细格算子 \(D\)、零空间 \(B\) 和聚合规则，能够在 CPU、GPU 或任意语言中构造同一 \(P/R\)、\(D_c=RDP\) 与 Schur 补系统。
      \end{block}
    \end{column}
    \begin{column}{0.31\textwidth}
      \begin{block}{控制流等价}
        伪代码明确写出 setup、batch、正交化、pre/post smooth、残差选择、递归、prepare/reconstruct 和停止条件。
      \end{block}
    \end{column}
    \begin{column}{0.31\textwidth}
      \begin{block}{实现可迁移}
        先给参考实现；再把 QUDA 的 checkerboard、halo、atomic、MMA、低精度和 cycle wrapper 映射回抽象接口。
      \end{block}
    \end{column}
  \end{columns}
  \vspace{0.08cm}
  \begin{alertblock}{边界声明}
    本报告确证源码控制流并给出可执行复现规范；本次未启动 CUDA/MPI，不虚报硬件收敛或性能。
  \end{alertblock}
  \vfill
  \begin{center}
    \begin{tikzpicture}[node distance=0.22cm and 0.28cm]
      \node[flowwide] (a) {输入与约定};
      \node[flowwide,right=of a] (b) {语言无关参考实现};
      \node[flowwide,right=of b] (c) {QUDA 优化映射};
      \node[flowwide,right=of c] (d) {逐层验证与回归};
      \draw[arrow] (a)--(b)--(c)--(d);
    \end{tikzpicture}
  \end{center}
  \source{本报告结构依据源码：\codepath{lib/multigrid.cpp:72-197, 1131-1224, 1275-1472}；图为复现路线抽象。}
\end{frame}

\begin{frame}[t]{先固定“算法合同”：只需要四类抽象接口}
  \small
  \begin{columns}[T,totalwidth=\textwidth]
    \begin{column}{0.48\textwidth}
      \begin{block}{数据与线性代数}
        \begin{tabularx}{\linewidth}{@{}lX@{}}
          \toprule
          接口 & 必须提供的语义 \\
          \midrule
          Field & 复数旋量/矩阵，支持 full 与 parity subset \\
          \(D\psi\) & 细格 residual operator，含边界条件与质量归一化 \\
          \(\langle u,v\rangle\) & 全局或局部复内积；范数为 \(\sum |u|^2\) \\
          \(P,R\) & 同一聚合基上的 prolong/restrict \\
          \bottomrule
        \end{tabularx}
      \end{block}
    \end{column}
    \begin{column}{0.48\textwidth}
      \begin{block}{实现层必须保持的不变量}
        \begin{enumerate}\setlength{\itemsep}{0pt}
          \item 粗点映射对每个细点唯一，且 inverse map 覆盖整个聚合；
          \item setup 中 (V) 的列在每个 aggregate 内正交归一；
          \item \(R=P^\dagger\)（本报告的 aggregate 路径）；
          \item coarse apply 与显式 (RDP) 的动作一致。
        \end{enumerate}
      \end{block}
      \begin{alertblock}{可移植策略}
        任何后端都可先用“逐粗基向量 apply \(P\to D\to R\)”构造 \(D_c\)，通过验证后再替换为 UV/VUV 融合核。
      \end{alertblock}
    \end{column}
  \end{columns}
  \source{Transfer 接口：\codepath{lib/transfer.cpp:259-328}；验证不变量：\codepath{lib/multigrid.cpp:785-880, 883-1068}。}
\end{frame}

\begin{frame}[t]{总体执行链不是黑盒：API、setup、solve、update 各自有边界}
  \centering
  \begin{tikzpicture}[node distance=0.28cm and 0.18cm]
    \node[flowwide,text width=2.0cm,font=\scriptsize] (g) {load gauge\\\(U,C,C^{-1}\)};
    \node[flowwide,text width=2.0cm,font=\scriptsize,right=of g] (api) {newMultigrid\\参数检查};
    \node[flowwide,text width=2.0cm,font=\scriptsize,right=of api] (setup) {每层 setup\\\(B,V,Y,X\)};
    \node[flowwide,text width=2.0cm,font=\scriptsize,right=of setup] (solve) {invert 一次/多次\\\(MG\) 应用};
    \node[decision,font=\scriptsize,right=of solve] (change) {场是否改变？};
    \node[flowwide,text width=2.0cm,font=\scriptsize,right=of change] (update) {thin/full\\update 或 destroy};
    \draw[arrow] (g)--(api)--(setup)--(solve)--(change)--node[above,font=\scriptsize]{是}(update);
    \draw[relation] (change.south) |- ++(0,-0.55) -| node[below,font=\scriptsize]{否：复用粗空间}(solve.south);
  \end{tikzpicture}
  \vspace{0.22cm}
  \begin{columns}[T,totalwidth=\textwidth]
    \begin{column}{0.47\textwidth}
      \begin{block}{构造时}
        公共入口创建细格 residual/smoother Dirac，再以 (B) 创建 (MG(0))；非粗层由 \texttt{\detokenize{reset()}} 创建 Transfer、粗场、粗 Dirac、下一层与 coarse solver。
      \end{block}
    \end{column}
    \begin{column}{0.47\textwidth}
      \begin{block}{更新时}
        thin update 只更新 Dirac 所引用的 gauge/clover 与质量；full update 重新创建细格 Dirac 并递归 \texttt{\detokenize{reset(refresh=true)}}，必要时刷新零空间。
      \end{block}
    \end{column}
  \end{columns}
  \source{公共入口：\codepath{lib/interface_quda.cpp:2853-2944}；更新：\codepath{lib/interface_quda.cpp:2946-3053}；递归 reset：\codepath{lib/multigrid.cpp:72-197}。}
\end{frame}

\begin{frame}[t]{每个非底层 \(MG_\ell\) 的状态：六个对象足以复现运行期}
  \begin{columns}[T,totalwidth=\textwidth]
    \begin{column}{0.55\textwidth}
      \centering
      \begin{tikzpicture}[node distance=0.20cm and 0.28cm]
        \node[state,text width=2.0cm] (b) {近零模\\\(B_\ell\)};
        \node[state,text width=2.0cm,right=of b] (t) {Transfer\\\(V,P,R\)};
        \node[state,text width=2.0cm,right=of t] (d) {粗算子\\\(Y,X\)};
        \node[state,text width=2.0cm,below=of d] (pc) {可选 PC\\\(\widehat Y,X^{-1}\)};
        \node[state,text width=2.0cm,left=of pc] (sm) {平滑器\\\(S_{pre},S_{post}\)};
        \node[state,text width=2.0cm,left=of sm] (c) {下一层\\\(MG_{\ell+1}\)};
        \draw[arrow] (b)--(t)--(d)--(pc);
        \draw[relation] (t)--(sm);
        \draw[relation] (d)--(sm);
        \draw[arrow] (sm)--(c);
        \draw[relation] (c.north) |- (b.south);
      \end{tikzpicture}
    \end{column}
    \begin{column}{0.41\textwidth}
      \small
      \begin{tabularx}{\linewidth}{@{}lX@{}}
        \toprule
        状态 & 用途 \\
        \midrule
        \(B_\ell\) & setup 的输入/输出，下一层的候选零空间 \\
        \(V_\ell\) & aggregate-local 的列正交基 \\
        \(Y,X\) & 粗 hopping 与 local term \\
        \(Yhat,X^{-1}\) & even/odd 预条件 dslash \\
        smoother & 负责局部高频误差 \\
        coarse solver & 递归 V-cycle 或底层 Krylov \\
        \bottomrule
      \end{tabularx}
    \end{column}
  \end{columns}
  \vspace{0.13cm}
  \begin{alertblock}{复现顺序}
    先实现 \(B\to V\to P/R\to D_c\)；再实现 \(X^{-1},Yhat\)；最后接入 smooth/recursive solve。每一步都能单独与显式矩阵动作对比。
  \end{alertblock}
  \source{成员与层级参数：\codepath{include/multigrid.h:75-217, 260-340}；创建顺序：\codepath{lib/multigrid.cpp:94-186}。}
\end{frame}

\begin{frame}[t]{维数协议（1/2）：先缩小时空，再重组自旋与颜色自由度}
  \small
  \begin{equation*}
    \begin{aligned}
      L_c^\mu&=L_f^\mu/b_\mu, & N_c^{\rm color}&=N_{\rm vec},\\
      N_c^{\rm spin}&=N_f^{\rm spin}/b_s &&(b_s>0),\\
      N_c^{\rm spin}&=2 &&(b_s=0\text{，staggered}).
    \end{aligned}
  \end{equation*}
  \begin{columns}[T,totalwidth=\textwidth]
    \begin{column}{0.48\textwidth}
      \begin{block}{场的抽象形状}
        细场 \(B_i\)：\(L_f^4\times N_f^s\times N_f^c\)；粗场：\(L_c^4\times N_c^s\times N_c^c\)。\(V\) 的每个 coarse color 是一列局部基。
      \end{block}
    \end{column}
    \begin{column}{0.48\textwidth}
      \begin{block}{存储边界}
        \(Y\) 是每个粗点的 4/8 个方向块，通常 nFace=1；\(X\) 是每个粗点的 local 方阵，无 ghost。checkerboard 只改变线性索引，不改变这些矩阵形状。
      \end{block}
    \end{column}
  \end{columns}
  \source{粗维数与 (Y/X) 创建：\codepath{lib/dirac_coarse.cpp:121-170}；布局解释为可迁移抽象。}
\end{frame}

\begin{frame}[t]{维数协议（2/2）：aggregate 容量是 \(N_{\rm vec}\) 的硬上限}
  \small
  \begin{equation*}
    C_{\rm agg}=\Bigl(\prod_{\mu=0}^{3}b_\mu\Bigr)N_f^{\rm color}
      \times\begin{cases}b_s,&b_s>0,\\[2pt]1/2,&b_s=0\ (\text{staggered}),\end{cases}
    \qquad N_{\rm vec}\le C_{\rm agg}.
  \end{equation*}
  \begin{columns}[T,totalwidth=\textwidth]
    \begin{column}{0.48\textwidth}
      \begin{block}{为什么是上限}
        \(C_{\rm agg}\) 是一个 aggregate 内可提供的独立局部自由度。超过它时，块 Gram--Schmidt 必然产生零列或线性相关列，后续 \(P/R\) 不可能形成单位局部基。
      \end{block}
    \end{column}
    \begin{column}{0.48\textwidth}
      \begin{alertblock}{先验检查}
        同时检查 \(1<\prod b_\mu\le1024\)、粗体积按 parity 可分和每维整除；不要只检查总 lattice volume。
      \end{alertblock}
    \end{column}
  \end{columns}
  \source{容量与块大小上限：\codepath{lib/transfer.cpp:25-83}、\codepath{lib/block_orthogonalize.in.cu:85-96}。}
\end{frame}

\begin{frame}[fragile,t]{奇偶与线性索引（1/3）：构造 fine-to-coarse map}
  \begin{equation*}
    p(x)=\left(\sum_{\mu=0}^{3}x_\mu\right)\bmod 2,
    \qquad x_c^\mu=\left\lfloor\frac{x_f^\mu}{b_\mu}\right\rfloor.
  \end{equation*}
  \begin{columns}[T,totalwidth=\textwidth]
    \begin{column}{0.39\textwidth}
      \begin{block}{正向 map}
        \texttt{\detokenize{fine_to_coarse[i]}} 把细格 parity-ordered 索引 \(i\) 映射到粗点 \(k_c\)。
      \end{block}
      \begin{alertblock}{必须记录}
        保留 \(x_f\)、\(p_f\)、\(x_c\) 三者，不要只存扁平 \(k_c\)；\(P/R\) 还要用原坐标恢复 neighbor。
      \end{alertblock}
    \end{column}
    \begin{column}{0.57\textwidth}
  \begin{lstlisting}[firstnumber=1]
for each fine linear index i:
    x_f, p_f = lattice_coordinate_and_parity(i)
    x_c = floor_div(x_f, geo_block_size)
    k_c = parity_ordered_index(x_c, parity(x_c))
    fine_to_coarse[i] = k_c
  \end{lstlisting}
    \end{column}
  \end{columns}
  \source{正向 map：\codepath{lib/transfer.cpp:200-241}。伪代码为等价抽象。}
\end{frame}

\begin{frame}[fragile,t]{奇偶与线性索引（2/3）：构造 coarse-to-fine 逆向列表}
  \begin{lstlisting}[firstnumber=1]
for each i:
    pairs.append((fine_to_coarse[i], i))
sort(pairs)                         # group by coarse site
coarse_to_fine = [pair.fine_index for pair in pairs]

for each coarse site q and parity p:
    assert count({i : fine_to_coarse[i] == q
                  and parity(i) == p}) > 0
  \end{lstlisting}
  \begin{columns}[T,totalwidth=\textwidth]
    \begin{column}{0.48\textwidth}
      \begin{block}{连续性}
        逆向表按 \((k_c,i_f)\) 排序，使同一 aggregate 的细点连续；GPU kernel 才能用一个 thread block 处理同一块。
      \end{block}
    \end{column}
    \begin{column}{0.48\textwidth}
      \begin{alertblock}{不变量}
        每个细点出现一次；每个粗点的列表覆盖所属细点，不能遗漏、重复或交换 parity。
      \end{alertblock}
    \end{column}
  \end{columns}
  \source{逆向 map：\codepath{lib/transfer.cpp:200-241}；kernel 读取：\codepath{include/kernels/coarse_op_kernel.cuh:1693-1704}。伪代码为等价抽象。}
\end{frame}

\begin{frame}[fragile,t]{奇偶与线性索引（3/3）：spin map 必须和 P/R 共用同一 parity}
  \begin{columns}[T,totalwidth=\textwidth]
    \begin{column}{0.48\textwidth}
\begin{lstlisting}[firstnumber=1]
if spin_block_size == 0:            # staggered
    spin_map[0][even] = 0
    spin_map[0][odd]  = 1
else:
    for s_f in fine_spin:
        for p in {even, odd}:
            spin_map[s_f][p] = floor_div(s_f, spin_block_size)
\end{lstlisting}
    \end{column}
    \begin{column}{0.48\textwidth}
      \begin{block}{读取时的索引链}
        先由 \(x_f\) 查 \(k_c\)，再由 \(p_f,s_f\) 查 \(s_c\)。同一个 \texttt{\detokenize{spin_map}} 必须被块正交化、\(P\)、\(R\) 和粗算子 setup 复用。
      \end{block}
      \begin{alertblock}{实现检查}
        对每个 coarse 点，\texttt{\detokenize{coarse_to_fine}} 必须列出属于它的细点；对每个 parity subset，spin map 不能切换到另一套编号。
      \end{alertblock}
    \end{column}
  \end{columns}
  \vspace{0.08cm}
  \begin{lstlisting}[firstnumber=1]
for each coarse site q and parity p:
    assert count({i : fine_to_coarse[i] == q and parity(i) == p}) > 0
    assert all(P/R use the same spin_map[s_f][p])
  \end{lstlisting}
  \source{spin map：\codepath{lib/transfer.cpp:243-257}；P/R parity subset：\codepath{lib/transfer.cpp:259-328}。伪代码为等价抽象。}
\end{frame}

\begin{frame}[fragile,t]{单层 setup 总览（1/4）：先固定几何与零空间}
  \small
  \begin{lstlisting}[firstnumber=1]
setup_level(level, D, U, C, B, cfg):
    validate_geometry_and_capacity(U.shape, cfg.geo_block_size,
                                   cfg.spin_block_size, cfg.n_vec)
    if cfg.compute_null:
        B = initialize_or_load_or_generate(B, cfg.null_source)
        B = generate_null_vectors(B, level, cfg)
  \end{lstlisting}
  \begin{columns}[T,totalwidth=\textwidth]
    \begin{column}{0.48\textwidth}
      \begin{block}{输入状态}\n+        本层接收细格算子 \(D\)、规范场/Cl​​over \(U,C\)、候选 \(B\) 和 level 配置；先验证几何，再决定 \(B\) 是读取、生成还是复用。\n+      \end{block}
    \end{column}
    \begin{column}{0.48\textwidth}
      \begin{alertblock}{顺序约束}\n+        \(N_{\rm vec}\)、block 和 spin map 未合法前不能分配 transfer；否则错误会延迟到 kernel 或粗解阶段。\n+      \end{alertblock}
    \end{column}
  \end{columns}
  \source{reset/零空间入口：\codepath{lib/multigrid.cpp:36-62, 72-94}；伪代码为语言无关等价实现。}
\end{frame}

\begin{frame}[fragile,t]{单层 setup 总览（2/4）：由 \(B\) 生成局部正交基与 P/R}
  \small
  \begin{lstlisting}[firstnumber=1]
    geo_map = build_geometry_maps(U.shape, cfg.geo_block_size)
    spin_map = build_spin_map(U.n_spin, cfg.spin_block_size)
    V = block_orthonormalize(B, geo_map, spin_map, cfg)
    P = make_prolongator(V, geo_map, spin_map)
    R = make_restrictor(V, geo_map, spin_map)       # R = P^dagger
    assert verify_maps(geo_map, spin_map)
    assert verify_local_orthogonality(V)
    return SetupState(B, V, P, R, geo_map, spin_map)
  \end{lstlisting}
  \begin{block}{此阶段的闭环}
    \(B\) 决定 \(V\)，\(V\) 决定 \(P/R\)。先通过 map 覆盖检查与 \(R=P^\dagger\)，再进入粗算子装配；所有后端必须复用同一 map。
  \end{block}
  \source{Transfer 构造：\codepath{lib/transfer.cpp:200-328}；块正交化：\codepath{include/kernels/block_orthogonalize.cuh:141-253}。伪代码为语言无关等价实现。}
\end{frame}

\begin{frame}[fragile,t]{单层 setup 总览（3/4）：由 transfer 装配粗算子}
  \small
  \begin{lstlisting}[firstnumber=1]
build_coarse_level(state, D, U, C, cfg):
    (Y, X) = build_galerkin_operator(U, C, state.V,
                                     state.geo_map, state.spin_map, cfg)
    verify_coarse_operator(state.P, state.R, D, Y, X, cfg)
    return CoarseState(Y, X, state.P, state.R)
  \end{lstlisting}
  \begin{columns}[T,totalwidth=\textwidth]
    \begin{column}{0.48\textwidth}
      \begin{block}{依赖链}
        \(P/R\) 确定 \(D_c=RDP\)；\(Y/X\) 是 coarse Dirac 的唯一 local/link 数据。验证应比较动作，而非只比较数组范数。
      \end{block}
    \end{column}
    \begin{column}{0.48\textwidth}
      \begin{alertblock}{可替换点}
        可先用逐粗基 \(P\to D\to R\) 构造 golden operator，再替换为 QUDA 的 UV/VUV 融合路径。
      \end{alertblock}
    \end{column}
  \end{columns}
  \source{粗算子创建：\codepath{lib/dirac_coarse.cpp:121-170, 225-308}；Galerkin：\codepath{lib/coarse_op.cuh:1043-1457}。}
\end{frame}

\begin{frame}[fragile,t]{单层 setup 总览（4/4）：PC 数据与下一层必须在粗算子之后}
  \small
  \begin{lstlisting}[firstnumber=1]
finish_level(coarse_state, cfg):
    Xinv = None; Yhat = None
    if cfg.needs_pc:
        Xinv = inverse_local_blocks(coarse_state.X)
        Yhat = build_preconditioned_links(coarse_state.Y, Xinv)
    coarse = allocate_next_level(coarse_geometry(),
                                 Nvec=cfg.n_vec, Nspin=coarse_spin)
    return Level(coarse_state, Xinv, Yhat, coarse)
  \end{lstlisting}
  \begin{alertblock}{不可交换的顺序}
    先完成 \(Y/X\) 和其 field metadata，再求 \(X^{-1}\)、\(\widehat Y\) 并创建下一层；否则 PC 的矩阵形状、halo 和 parity 信息没有来源。
  \end{alertblock}
  \source{递归 reset/下一层：\codepath{lib/multigrid.cpp:94-186}；PC 构造：\codepath{lib/dirac_coarse.cpp:225-308}。伪代码为组合实现。}
\end{frame}

\begin{frame}[fragile,t]{几何检查（1/2）：非法 block 的逐维减半行为}
  \begin{columns}[T,totalwidth=\textwidth]
    \begin{column}{0.55\textwidth}
      \begin{lstlisting}[firstnumber=1]
for d in dimensions 0..3:
    while b[d] > 0:
        coarse_d = L_f[d] / b[d]
        invalid = (d == 0 and L_f[0] == b[0])
               or ((coarse_d + 1) % 2 == 0)
               or (coarse_d * b[d] != L_f[d])
        if not invalid:
            break
        warn_or_record(d, b[d])
        b[d] = b[d] / 2
    if b[d] == 0:
        fail("dimension cannot be blocked")
      \end{lstlisting}
    \end{column}
    \begin{column}{0.41\textwidth}
      \begin{block}{判定含义}
        每一维都要求整除；粗维还必须能按 parity 分开。遇到非法值时记录 warning 并减半，直到可用或失败。
      \end{block}
      \begin{alertblock}{两种复现目标}
        数学复现可直接拒绝非法配置；行为复现则必须保留逐维减半、warning 和最终 block。
      \end{alertblock}
    \end{column}
  \end{columns}
  \source{检查与减半：\codepath{lib/transfer.cpp:25-83}；粗维 parity 约束：\codepath{lib/transfer.cpp:45-83}。伪代码显式保留行为分支。}
\end{frame}

\begin{frame}[fragile,t]{几何检查（2/2）：aggregate 容量决定 \(N_{\rm vec}\) 能否存在}
  \begin{lstlisting}[firstnumber=1]
aggregate = product(b)
if spin_block == 0:                    # staggered
    capacity = aggregate * N_f_color / 2
else:
    capacity = aggregate * N_f_color * spin_block

assert Nvec <= capacity
assert 1 < aggregate <= 1024
assert aggregate % 2 == 0
  \end{lstlisting}
  \vspace{0.10cm}
  \begin{columns}[T,totalwidth=\textwidth]
    \begin{column}{0.48\textwidth}
      \begin{block}{物理/线性代数含义}
        一个 aggregate 能提供的独立局部自由度就是粗基列数的上限；若 \(N_{\rm vec}>C_{\rm agg}\)，局部 Gram--Schmidt 必然出现零向量。
      \end{block}
    \end{column}
    \begin{column}{0.48\textwidth}
      \begin{alertblock}{实现含义}
        \(C_{\rm agg}\) 与 \(b_\mu\)、细颜色、spin block 共同决定粗场维数；不要只检查总 lattice volume。
      \end{alertblock}
    \end{column}
  \end{columns}
  \source{容量与块大小上限：\codepath{lib/transfer.cpp:25-83}、\codepath{lib/block_orthogonalize.in.cu:85-96}。}
\end{frame}

\begin{frame}[fragile,t]{近零模生成（1/3）：先确定 setup solver 的数值语义}
  \begin{lstlisting}[firstnumber=1]
configure_null_setup(level, cfg, refresh=false):
    sp.maxiter = refresh ? cfg.setup_maxiter_refresh[level]
                         : cfg.setup_maxiter[level]
    sp.tol = cfg.setup_tol[level]
    sp.use_init_guess = YES
    sp.delta = 1e-1
    sp.inv_type = cfg.setup_inv_type[level]
    sp.precision = cfg.cuda_prec_sloppy
    sp.residual_type = L2_RELATIVE
    sp.compute_null_vector = YES
    return sp
  \end{lstlisting}
  \vspace{0.08cm}
  \begin{columns}[T,totalwidth=\textwidth]
    \begin{column}{0.48\textwidth}
      \begin{block}{为什么是“近零模”}
        setup solver 反复削弱 \(D\) 的高频分量，使 \(B_i\) 落入低频子空间；它不是任意随机向量的存档。
      \end{block}
    \end{column}
    \begin{column}{0.48\textwidth}
      \begin{alertblock}{必须保持的状态}
        \(use\_init\_guess=YES\)、相对残差类型、sloppy 精度和 refresh 的独立 maxiter 都会改变最终 \(B\)，不能只复现 solver 名称。
      \end{alertblock}
    \end{column}
  \end{columns}
  \source{setup solver 参数：\codepath{lib/multigrid.cpp:1275-1363}；参数来源：\codepath{tests/utils/set_params.cpp:516-582}。伪代码为行为抽象。}
\end{frame}

\begin{frame}[fragile,t]{近零模生成（2/3）：分批 solve 并把结果写回同一组 \(B_i\)}
  \begin{lstlisting}[firstnumber=1]
generate_null_vectors(B, level, cfg, refresh=false):
    sp = configure_null_setup(level, cfg, refresh)
    batch = cfg.n_vec_batch[level]
    require len(B) % batch == 0

    for setup_iter in 1 .. cfg.num_setup_iter[level]:
        if cfg.pre_orthonormalize:
            global_gram_schmidt(B)
        for first in range(0, len(B), batch):
            B_i = B[first:first+batch]
            if cfg.setup_type == TEST_VECTOR_SETUP:
                rhs = copy(B_i); x = zero_like(B_i)
            else:
                x = copy(B_i); rhs = zero_like(B_i)
            solve_setup_system(x, rhs, sp)
            B[first:first+batch] = x
        if cfg.post_orthonormalize:
            global_gram_schmidt(B)
        refresh_transfer_and_next_level_if_needed(B, cfg)
    return B
  \end{lstlisting}
  \vspace{0.08cm}
  \begin{block}{两个 setup 分支的语义}
    \texttt{\detokenize{TEST_VECTOR_SETUP}} 以 \(B_i\) 为右端、零初值求解；普通分支以 \(B_i\) 为初值、零右端进行平滑。两者都把 solver 输出写回 \(B_i\)，所以批次切分不能改变向量顺序。
  \end{block}
  \begin{alertblock}{传播时机}
    每次 setup iteration 后才进行 post-orthonormalize 与下一层刷新；若提前传播，下一层看到的不是同一个 \(B\)。
  \end{alertblock}
  \source{batch 循环与写回：\codepath{lib/multigrid.cpp:1365-1413}；setup type：\codepath{tests/utils/set_params.cpp:516-582}。}
\end{frame}

\begin{frame}[t]{近零模生成（2/2）：CG、GCR、MG 三条内部求解路径不能混写}
  \begin{columns}[T,totalwidth=\textwidth]
    \begin{column}{0.47\textwidth}
      \begin{block}{CG / CA-CG}
        setup solver 不是直接作用于 \(D\)，而是包装 \texttt{\detokenize{DiracMdagM}}，即以 \(D^\dagger D\) 为求解矩阵。CA-CG 另外设置 Chebyshev/CA basis 参数。
      \end{block}
      \begin{block}{普通 Krylov}
        GCR、BiCGStab-L 等由 \texttt{\detokenize{Solver::create}} 直接以 \texttt{\detokenize{matSmooth}} 为矩阵；若输入配置带 Schwarz，源码在此阶段先暂时关闭它。
      \end{block}
    \end{column}
    \begin{column}{0.49\textwidth}
      \begin{block}{MG setup 自举}
        当 \texttt{\detokenize{setup_inv_type=MG}}：
        \[
          \text{setup solver}=\text{GCR},
          \qquad K=MG_\ell,
          \qquad \text{每轮最多一次 MG 预条件调用}.
        \]
        先保证本层 smoother 存在，再用本层 MG 作为 GCR 的 preconditioner。
      \end{block}
      \begin{alertblock}{递归更新}
        \texttt{\detokenize{generate_all_levels=true}}：向下层重新生成；否则先 \(B_{\ell+1}=R_\ell B_\ell\)，再重建下一层 Transfer。自举完成后恢复 Schwarz/halo 精度并释放临时 solver。
      \end{alertblock}
    \end{column}
  \end{columns}
  \source{solver 分支：\codepath{lib/multigrid.cpp:1335-1363}；MG 自举与向下传播：\codepath{lib/multigrid.cpp:1429-1466}；Schwarz 临时关闭：\codepath{lib/multigrid.cpp:1321-1333}。}
\end{frame}

\begin{frame}[t]{全局正交化与块正交化是两层不同约束}
  \begin{columns}[T,totalwidth=\textwidth]
    \begin{column}{0.46\textwidth}
      \begin{block}{全局正交化（可选）}
        作用域是整个 local/global lattice 的 \(B_i\)：
        \[
          B_i\leftarrow B_i-\sum_{j<i}\langle B_j,B_i\rangle B_j,
          \quad B_i\leftarrow B_i/\|B_i\|.
        \]
        由 \texttt{\detokenize{pre_orthonormalize}}、\texttt{\detokenize{post_orthonormalize}} 控制，内积可能跨 MPI reduction。
      \end{block}
    \end{column}
    \begin{column}{0.50\textwidth}
      \begin{block}{块正交化（Transfer 构造必需）}
        每个 aggregate、每个 coarse-spin/chirality block 单独执行 Gram--Schmidt；它保证的是 \(V^\dagger V\approx I\) 的局部关系，而不是不同 aggregate 之间的全局正交。
      \end{block}
      \begin{center}
        \begin{tikzpicture}[node distance=0.15cm and 0.19cm]
          \node[phase] (g) {全局 \(B_i\)：\(\langle B_j,B_i\rangle\)};
          \node[phase,right=of g] (a) {按 aggregate 聚类细点};
          \node[phase,right=of a] (b) {局部 \(V_i\)：内积与归一};
          \draw[arrow] (g)--(a)--(b);
        \end{tikzpicture}
      \end{center}
    \end{column}
  \end{columns}
  \vspace{0.16cm}
  \begin{alertblock}{不要用全局 \(B^\dagger B=I\) 替代块正交}
    Galerkin transfer 的局部列基按 aggregate 存储；只做全局正交并不能使同一 aggregate 内的列成为 \(P^\dagger P\) 的单位块。
  \end{alertblock}
  \source{全局正交化：\codepath{lib/multigrid.cpp:1369-1381, 1415-1427}；块正交化：\codepath{include/kernels/block_orthogonalize.cuh:141-253}。}
\end{frame}

\begin{frame}[fragile,t]{块正交化伪代码（1/3）：按 aggregate 和向量批次装载候选列}
  \begin{lstlisting}[firstnumber=1]
block_orthonormalize(B, geo_map, spin_map, n_block_ortho):
    V = allocate_matrix_field(n_vec, fine_spin, fine_color)
    for each coarse site q:
      for each coarse chirality h:
        for repeat in 0 .. n_block_ortho-1:
          for j0 in 0, mVec, 2*mVec, ... n_vec-mVec:
            v[0:mVec] = load(B if repeat == 0 else V,
                              q, aggregate(q), h, j0:j0+mVec)
            for i in 0 .. j0-1:
              vi = load(V, q, aggregate(q), h, i)
              for m in 0 .. mVec-1:
                coeff[m] = sum_{x in aggregate(q)}
                           inner(vi(x), v[m](x))
              coeff = reduce_within_aggregate(coeff)
              for m in 0 .. mVec-1:
                v[m] = v[m] - coeff[m] * vi
            local_block_orthogonalize_and_normalize(v)
            store(V, q, aggregate(q), h, j0:j0+mVec, v)
  \end{lstlisting}
  \vspace{0.08cm}
  \begin{block}{\texttt{\detokenize{mVec}} 的意义}
    \(mVec\) 只是批处理宽度；它改变 tile 并行粒度，不改变“先对旧列正交、再写回”的 Gram--Schmidt 顺序。先令 \texttt{\detokenize{mVec=1}} 可得到最容易核对的 reference。
  \end{block}
  \source{批处理与首次从 \(B\) 加载：\codepath{include/kernels/block_orthogonalize.cuh:170-190}；前向投影：\codepath{include/kernels/block_orthogonalize.cuh:192-213}。伪代码为等价抽象。}
\end{frame}

\begin{frame}[fragile,t]{块正交化伪代码（2/3）：tile 内的列间正交和归一化}
  \begin{lstlisting}[firstnumber=1]
local_block_orthogonalize_and_normalize(v[0:mVec]):
    for m in 0 .. mVec-1:
        for i in 0 .. m-1:
            c = reduce_{x in aggregate}(inner(v[i](x), v[m](x)))
            v[m] = v[m] - c * v[i]
        n2 = reduce_{x in aggregate}(norm2(v[m](x)))
        if n2 <= 0: fail("zero null vector in aggregate")
        v[m] = v[m] / sqrt(n2)
  \end{lstlisting}
  \vspace{0.10cm}
  \begin{columns}[T,totalwidth=\textwidth]
    \begin{column}{0.48\textwidth}
      \begin{block}{两类正交项}
        外层 \(i<j_0\) 处理当前 tile 与历史列；tile 内 \(i<m\) 处理同批列之间的交叉项。少掉任一层都会使 \(V^\dagger V\) 偏离单位阵。
      \end{block}
    \end{column}
    \begin{column}{0.48\textwidth}
      \begin{alertblock}{归一化边界}
        \(n^2\) 是 aggregate 内的全局（或 rank-local 后再 reduction）范数；\(n^2\le0\) 必须报告为零基向量，不能静默除零。
      \end{alertblock}
    \end{column}
  \end{columns}
  \source{tile 内正交与归一化：\codepath{include/kernels/block_orthogonalize.cuh:215-251}。}
\end{frame}

\begin{frame}[fragile,t]{块正交化伪代码（3/3）：重复正交与低精度 two-pass}
  \begin{lstlisting}[firstnumber=1]
two_pass_block_ortho(B, cfg):
    V = block_orthonormalize(B, ..., cfg.n_block_ortho)
    if cfg.block_ortho_two_pass and V.precision < SINGLE:
        scale = abs_max(V)
        if 1.05 * scale < 1:
            V.set_storage_scale(1.05 * scale)
            V = block_orthonormalize(B, ..., cfg.n_block_ortho)
    return V
  \end{lstlisting}
  \vspace{0.10cm}
  \begin{block}{重复次数}
    \texttt{\detokenize{n_block_ortho}} 重复整套块 Gram--Schmidt；它抑制舍入残差，但不是改变数学定义的额外基。
  \end{block}
  \begin{alertblock}{two-pass 的边界}
    先用实际最大幅值设置 fixed-point/half scale，再从原始 \(B\) 重算；这是存储精度补救，不是额外的收敛判据，scale 必须随 \(V\) 保存。
  \end{alertblock}
  \source{two-pass scale：\codepath{lib/block_orthogonalize.in.cu:104-130}；重复正交：\codepath{include/kernels/block_orthogonalize.cuh:141-253}。}
\end{frame}

\begin{frame}[t]{P/R 的数学定义（1/2）：prolongation 是按局部基展开}
  \begin{equation*}
    (P\psi_c)_{s_f c_f}(x_f)
    =\sum_{a=0}^{N_c^c-1}
    V_{s_f c_f,a}(x_f)\,
    \psi_{c,\,s_c(s_f,p_f),a}(x_c(x_f)).
  \end{equation*}
  \vspace{0.10cm}
  \begin{columns}[T,totalwidth=\textwidth]
    \begin{column}{0.48\textwidth}
      \begin{block}{局部基的角色}
        \(V(x_f)\) 的列把 coarse color \(a\) 旋转到细 spin/color；\(x_c(x_f)\) 和 \(s_c(s_f,p_f)\) 只负责位置与 coarse-spin 映射。
      \end{block}
    \end{column}
    \begin{column}{0.48\textwidth}
      \begin{alertblock}{单 parity}
        full field 计算两个 parity；single-parity 由 \texttt{\detokenize{Transfer::setSiteSubset}} 固定目标 parity，不能把单 parity 的 \(V\) 当作 full-field 基。
      \end{alertblock}
    \end{column}
  \end{columns}
  \source{P 入口：\codepath{lib/transfer.cpp:259-292}；kernel 旋转：\codepath{include/kernels/prolongator.cuh:65-132}。}
\end{frame}

\begin{frame}[t]{P/R 的数学定义（2/2）：restriction 是共轭转置和 aggregate reduction}
  \begin{equation*}
    (R\psi_f)_{s_c a}(x_c)
    =\sum_{x_f\in A(x_c)}\sum_{s_f,c_f}
    V^*_{s_f c_f,a}(x_f)\,
    \psi_{f,s_f c_f}(x_f).
  \end{equation*}
  \vspace{0.08cm}
  \begin{block}{局部单位性}
    若每个 aggregate 内 \(V^\dagger V=I\)，则
    \[
      (RP\eta_c)(x_c)=\eta_c(x_c).
    \]
    这是 coarse span 检查的理论基准；它同时要求 map、清零、累加和共轭方向全部正确。
  \end{block}
  \begin{center}
    \begin{tikzpicture}[node distance=0.20cm and 0.20cm]
      \node[flowwide,text width=1.75cm,font=\scriptsize] (c) {粗场 \(\eta_c\)};
      \node[state,text width=1.90cm,font=\scriptsize,right=of c] (map) {查 map 与 spin map};
      \node[state,text width=1.90cm,font=\scriptsize,right=of map] (rot) {乘 \(V\) 或 \(V^\dagger\)};
      \node[flowwide,text width=1.75cm,font=\scriptsize,right=of rot] (f) {细场 \(\eta_f\)};
      \draw[arrow] (c)--node[above,font=\scriptsize]{\(P\)}(map)--(rot)--(f);
      \draw[relation] (f.south) |- ++(0,-0.52) -| node[below,font=\scriptsize]{aggregate reduction：\(R\)} (c.south);
    \end{tikzpicture}
  \end{center}
  \source{R 入口：\codepath{lib/transfer.cpp:292-328}；kernel reduction：\codepath{include/kernels/restrictor.cuh:87-150}。}
\end{frame}

\begin{frame}[fragile,t]{P 参考实现：先保证位置映射、清零和 parity 完全正确}
  \begin{lstlisting}[firstnumber=1]
prolongate(out_f, in_c):
    for each fine site x_f in requested parity subset:
        q = geo_map[x_f]; p = parity(x_f)
        for s_f, c_f:
            s_c = spin_map[s_f][p]
            out_f[x_f,s_f,c_f] = 0
            for a in coarse_colors:
                out_f[x_f,s_f,c_f] += V[x_f,s_f,c_f,a] \
                                      * in_c[q,s_c,a]
  \end{lstlisting}
  \vspace{0.10cm}
  \begin{columns}[T,totalwidth=\textwidth]
    \begin{column}{0.48\textwidth}
      \begin{block}{GPU 映射}
        一个线程可负责细点/粗色块；一个 block 可负责一个 aggregate；MMA 只替换内层矩阵乘。
      \end{block}
    \end{column}
    \begin{column}{0.48\textwidth}
      \begin{alertblock}{不变量}
        每个细点先清零再对 coarse color \(a\) 求和；\(q,p,s_c\) 必须来自同一组 map。
      \end{alertblock}
    \end{column}
  \end{columns}
  \source{Prolongate：\codepath{include/kernels/prolongator.cuh:65-132}。伪代码为等价实现。}
\end{frame}

\begin{frame}[fragile,t]{R 参考实现：清零后做共轭基累加和 reduction}
  \begin{lstlisting}[firstnumber=1]
restrict(out_c, in_f):
    fill(out_c, 0)
    for each fine site x_f in requested parity subset:
        q = geo_map[x_f]; p = parity(x_f)
        for s_f, c_f, a:
            s_c = spin_map[s_f][p]
            out_c[q,s_c,a] += conj(V[x_f,s_f,c_f,a]) \
                              * in_f[x_f,s_f,c_f]
    reduce_out_c_across_threads_and_ranks(out_c)
    return out_c
  \end{lstlisting}
  \vspace{0.10cm}
  \begin{alertblock}{第一条单元测试}
    随机粗场检查 \(\|R(P\eta_c)-\eta_c\|/\|\eta_c\|\)；随机细场检查 \(R\) 与显式共轭转置的差异。若输出未清零，误差会随重复调用累积。
  \end{alertblock}
  \source{Restrict 的 zero/reduction：\codepath{lib/transfer.cpp:292-328}、\codepath{include/kernels/restrictor.cuh:126-150}。伪代码为等价实现。}
\end{frame}

\begin{frame}[fragile,t]{Galerkin golden reference（1/2）：逐粗基列执行 \(P\to D\to R\)}
  \begin{lstlisting}[firstnumber=1]
build_galerkin_reference(D, P, R, coarse_shape):
    Dc = zero_matrix_or_operator(coarse_shape)
    for each coarse source basis index beta=(q,s_c,a):
        e_c = unit_vector(beta)
        e_f = P(e_c)
        y_f = D(e_f)
        column = R(y_f)
        Dc[:, beta] = column
    return Dc
  \end{lstlisting}
  \vspace{0.10cm}
  \begin{columns}[T,totalwidth=\textwidth]
    \begin{column}{0.48\textwidth}
      \begin{block}{为什么可迁移}
        只依赖 \(D,P,R\) 的动作，不依赖 CUDA thread/block、GaugeField order 或 atomic；复杂度大，但可作为 CPU golden reference。
      \end{block}
    \end{column}
    \begin{column}{0.48\textwidth}
      \begin{alertblock}{列/行约定}
        \(e_c\) 是 coarse 的单位列；\(R(D(Pe_c))\) 直接写入 \(D_c\) 的同一列。若转置或共轭方向错，\(D_c-RDP\) 会立即暴露。
      \end{alertblock}
    \end{column}
  \end{columns}
  \source{Galerkin verify 的比较对象：\codepath{lib/multigrid.cpp:883-928}。伪代码定义了可移植 golden path。}
\end{frame}

\begin{frame}[fragile,t]{Galerkin golden reference（2/2）：把显式 \(D_c\) 拆成 \(Y\) 与 \(X\)}
  \begin{lstlisting}[firstnumber=1]
split_into_Y_and_X(Dc):
    for each output q and input q2:
        if q2 == q:
            X[q] += Dc[q,q2]
        elif q2 == q + unit_direction(mu):
            Y_fwd[mu,q] = Dc[q,q2]
        elif q2 == q - unit_direction(mu):
            Y_back[mu,q] = Dc[q,q2]
        else:
            preserve_explicit_long_range_term_or_fail()
    return Y_fwd, Y_back, X
  \end{lstlisting}
  \vspace{0.10cm}
  \begin{equation*}
    D_c=RDP,\qquad
    Y_{q\to q'}=\sum_{x\in A_q,y\in A_{q'}}V^\dagger(x)H(x,y)V(y),
    \qquad
    X_q=\sum_{x,y\in A_q}V^\dagger(x)D(x,y)V(y).
  \end{equation*}
  \begin{alertblock}{结构约束}
    标准 Wilson/Clover aggregate 把跨 coarse-site 的一跳项放进 \(Y\)，同 aggregate 内项放进 \(X\)。含长链时必须显式决定保留长链接，或按 QUDA 的 truncation 规则拒绝/截断。
  \end{alertblock}
  \source{粗场拆分与 \texttt{\detokenize{isCoarseDiagonal}}：\codepath{include/kernels/coarse_op_kernel.cuh:1608-1647}；长链接条件：\codepath{lib/coarse_op.cuh:1193-1222}。}
\end{frame}

\begin{frame}[t]{QUDA 的快速 Galerkin 图（1/2）：UV/VUV 是 golden path 的融合与复用}
  \centering
  \begin{tikzpicture}[node distance=0.20cm and 0.12cm]
    \node[phase,text width=1.92cm] (v) {邻点 \(V\) / ghost exchange};
    \node[phase,right=of v] (uv) {UV\(U\cdot V_{\rm nb}\)};
    \node[phase,right=of uv] (av) {AV\(C V\) 或\(X^{-1}V\)};
    \node[phase,right=of av] (vuv) {VUV\(V^\dagger\,UV\)};
    \node[phase,right=of vuv] (split) {按 coarse 邻接累加 \(Y/X\)};
    \node[phase,below=0.34cm of split,text width=2.20cm] (extras) {identity / mass / twist / clover};
    \node[phase,left=of extras,text width=2.20cm] (rev) {必要时 reverse 补 forward link};
    \draw[arrow] (v)--(uv)--(av)--(vuv)--(split);
    \draw[arrow] (split)--(extras);
    \draw[arrow] (extras)--(rev);
  \end{tikzpicture}
  \vspace{0.16cm}
  \begin{block}{正确性关系}
    UV/VUV 只把 \(RDP\) 中可复用的邻点乘法和局部矩阵乘法融合；它必须与 golden reference 的 \(D_c\) 在每个 coarse matrix element 上一致。
  \end{block}
  \source{UV/VUV 调度：\codepath{lib/coarse_op.cuh:1043-1062, 1166-1215, 1269-1323}；VUV kernel：\codepath{include/kernels/coarse_op_kernel.cuh:1067-1138}。}
\end{frame}

\begin{frame}[t]{QUDA 的快速 Galerkin 图（2/2）：projector、系数与存储边界}
  \begin{equation*}
    D_c=RDP,\qquad
    H_{+\mu}=(1-\gamma_\mu)U_\mu,\qquad
    H_{-\mu}=(1+\gamma_\mu)U_\mu^\dagger .
  \end{equation*}
  \begin{columns}[T,totalwidth=\textwidth]
    \begin{column}{0.47\textwidth}
      \begin{block}{Wilson/Clover}
        forward 使用 \(1-\gamma_\mu\)，backward 使用 \(1+\gamma_\mu\)；spin 对角和非对角项分别累加到 coarse spin block。
      \end{block}
    \end{column}
    \begin{column}{0.47\textwidth}
      \begin{alertblock}{\(\kappa\) 与 \(X\)}
        外部 \(Y\) 通常不含应用时的 \(-\kappa\)；粗 dslash 汇总后再乘 \(-\kappa\)。内部项已在写入 \(X\) 前乘相应系数。
      \end{alertblock}
    \end{column}
  \end{columns}
  \source{Wilson projector：\codepath{include/kernels/coarse_op_kernel.cuh:1067-1138}；应用时 \(-\kappa\)：\codepath{include/kernels/dslash_coarse.cuh:292-324}。}
\end{frame}

\begin{frame}[t]{Galerkin 的关键分类（1/2）：fine hopping 决定写 \(Y\) 还是 \(X\)}
  \begin{equation*}
    H_{+\mu}(x,y)=(1-\gamma_\mu)U_\mu(x)\,\delta_{y,x+\hat\mu},
    \qquad
    H_{-\mu}(x,y)=(1+\gamma_\mu)U_\mu^\dagger(x-\hat\mu)\,
    \delta_{y,x-\hat\mu}.
  \end{equation*}
  \vspace{0.12cm}
  \begin{columns}[T,totalwidth=\textwidth]
    \begin{column}{0.48\textwidth}
      \begin{block}{同一粗点}
        若 \(\lfloor x/b\rfloor=\lfloor (x\pm\hat\mu)/b\rfloor\)，该项是 local term，累加 \(X(q)\)，并保留完整 spin/color 矩阵。
      \end{block}
    \end{column}
    \begin{column}{0.48\textwidth}
      \begin{alertblock}{相邻粗点}
        若终点落入 \(q\pm\hat\mu\)，该项写入对应方向的 \(Y\)；apply 时从 neighbor coarse spinor gather。
      \end{alertblock}
    \end{column}
  \end{columns}
  \vspace{0.18cm}
  \begin{center}
    \begin{tikzpicture}[node distance=0.38cm and 0.50cm]
      \node[flow,text width=1.65cm,font=\scriptsize] (a) {细点 \(x\in A_q\)};
      \node[flow,text width=1.80cm,font=\scriptsize,right=of a] (b) {hop 到 \(y=x\pm\hat\mu\)};
      \node[decision,font=\scriptsize,right=of b] (c) {same \(q\)?};
      \node[state,text width=1.55cm,font=\scriptsize,below right=0.22cm and 0.28cm of c] (x) {写入 \(X(q)\)};
      \node[state,text width=1.55cm,font=\scriptsize,above right=0.22cm and 0.28cm of c] (y) {写入 \(Y_{\pm\mu}\)};
      \draw[arrow] (a)--(b)--(c);
      \draw[arrow] (c)--node[above,font=\scriptsize]{否}(y);
      \draw[arrow] (c)--node[below,font=\scriptsize]{是}(x);
    \end{tikzpicture}
  \end{center}
  \source{对角判定与写回：\codepath{include/kernels/coarse_op_kernel.cuh:1548-1647}。图为同一判定的可执行解释。}
\end{frame}

\begin{frame}[t]{Galerkin 的关键分类（2/2）：长链接与方向存储不可含糊}
  \begin{block}{一跳 Wilson/Clover}
    \(Y_{+\mu}(q)\) 表示从 \(q\) gather \(q+\hat\mu\) 的 forward block；\(Y_{-\mu}(q)\) 表示从 \(q\) gather \(q-\hat\mu\) 的 backward block。方向、坐标和 dagger 必须和 coarse dslash 的 gather 完全一致。
  \end{block}
  \vspace{0.08cm}
  \begin{columns}[T,totalwidth=\textwidth]
    \begin{column}{0.48\textwidth}
      \begin{block}{长链接}
        ASQTAD/HISQ 需要 \(nFace=3\)。若聚合方向小于 3，不能假设一跳 \(Y\) 足够；按 \texttt{\detokenize{allow_truncation}} 选择截断、保留长链接或拒绝配置。
      \end{block}
    \end{column}
    \begin{column}{0.48\textwidth}
      \begin{alertblock}{互斥分类}
        每个 fine hopping 恰好进入 \(X\)、某个 \(Y_{\pm\mu}\) 或明确的 long-range bucket；重复计数和静默丢项都会破坏 \(D_c-RDP\)。
      \end{alertblock}
    \end{column}
  \end{columns}
  \vspace{0.15cm}
  \begin{equation*}
    D_c(q,q')=
    \begin{cases}
      X(q),&q'=q,\\
      Y_{+\mu}(q),&q'=q+\hat\mu,\\
      Y_{-\mu}(q),&q'=q-\hat\mu.
    \end{cases}
  \end{equation*}
  \source{长链接条件：\codepath{lib/coarse_op.cuh:1193-1222, 1295-1323}；方向与 coarse dslash：\codepath{include/kernels/dslash_coarse.cuh:126-324}。}
\end{frame}

\begin{frame}[t]{\(X\) 的组成（1/2）：内部 hopping 与 clover 投影}
  \begin{equation*}
    X(q)=X_{\rm internal\ hop}(q)+X_{\rm clover}(q)+X_{\rm local\ diag}(q).
  \end{equation*}
  \vspace{0.10cm}
  \begin{tabularx}{\textwidth}{@{}lX@{}}
    \toprule
    分量 & QUDA 逻辑与不可重复计算的边界 \\
    \midrule
    internal hop & \(V^\dagger H V\)，写入 \(X\)，已带 hopping 系数；只接收 aggregate 内部项 \\
    clover & fine Clover 的 \(V^\dagger C V\)；仅在选择 clover 构造时加入 \\
    \bottomrule
  \end{tabularx}
  \vspace{0.12cm}
  \begin{block}{与 \(Y\) 的互斥性}
    终点仍在 \(A_q\) 的 fine hopping 才进入 \(X(q)\)；终点落入邻 aggregate 的项必须进入 \(Y_{\pm\mu}\)，不能在两个位置重复累加。
  \end{block}
  \source{internal/clover：\codepath{lib/coarse_op.cuh:1383-1441}；对角/非对角分支：\codepath{include/kernels/coarse_op_kernel.cuh:1548-1647}。}
\end{frame}

\begin{frame}[t]{\(X\) 的组成（2/2）：identity、mass、twisted-mass 逐项补齐}
  \begin{equation*}
    X_{\rm local\ diag}(q)
    =I_{\rm id}+2m\,I_{\rm stag}
      +i\mu\,{\rm diag}(+1_{\rm upper},-1_{\rm lower}).
  \end{equation*}
  \vspace{0.10cm}
  \begin{tabularx}{\textwidth}{@{}lX@{}}
    \toprule
    分量 & 触发条件与实现语义 \\
    \midrule
    identity & 非 staggered 路径显式加 coarse identity \\
    staggered mass & staggered 路径对角加 \(2mI\) \\
    twisted mass & coarse spin 上半块加 \(+i\mu\)，下半块加 \(-i\mu\) \\
    \bottomrule
  \end{tabularx}
  \vspace{0.12cm}
  \begin{alertblock}{装配边界}
    \texttt{COMPUTE\_COARSE\_CLOVER} 只构造 fine clover 投影；internal hopping 已在 VUV diagonal 分支进入 \(X\)，不能再次加。
  \end{alertblock}
  \source{identity/mass/twist：\codepath{lib/coarse_op.cuh:1383-1441}；对应 kernel：\codepath{include/kernels/coarse_op_kernel.cuh:1768-1950}。}
\end{frame}

\begin{frame}[fragile,t]{\(X\) 的装配伪代码（1/2）：逐 fine hopping 判断是否留在 aggregate}
\begin{lstlisting}[firstnumber=1]
build_local_hop_block(V, fine_operator, aggregate q):
    X = zero_matrix(coarse_dof, coarse_dof)
    for x in sites(q):
        for mu in 0 .. 3:
            y = neighbor(x, +mu)
            if aggregate_id(y) == q:
                H = forward_hop(fine_operator, x, mu)
                X += dagger(V[x]) * H * V[y]
            y = neighbor(x, -mu)
            if aggregate_id(y) == q:
                H = backward_hop(fine_operator, x, mu)
                X += dagger(V[x]) * H * V[y]
    return X
\end{lstlisting}
  \vspace{0.12cm}
  \begin{alertblock}{关键不变量}
    此函数只接收 aggregate 内部 hopping；跨 aggregate 的项由同一套 \(V^\dagger H V\) 累加到 \(Y_{\pm\mu}\)。
  \end{alertblock}
  \source{对角/非对角分支：\codepath{include/kernels/coarse_op_kernel.cuh:1548-1647}。伪代码为可移植等价实现。}
\end{frame}

\begin{frame}[fragile,t]{\(X\) 的装配伪代码（2/2）：在 hopping 后按算子类型补 local term}
\begin{lstlisting}[firstnumber=1]
finish_local_block(X, V, fine_operator, aggregate q, cfg):
    if cfg.has_clover:
        for x in sites(q):
            X += dagger(V[x]) * clover(x) * V[x]
    else if not cfg.is_staggered:
        X += identity_matrix()
    if cfg.is_staggered:
        X += 2 * cfg.mass * identity_matrix()
    if cfg.is_twisted:
        add_twisted_spin_diagonal(X, cfg.mu, cfg.spin_sign)
    return X
\end{lstlisting}
  \vspace{0.12cm}
  \begin{block}{执行顺序}
    先完成 \(V^\dagger H V\) 的 internal 分支，再执行 clover/identity/mass/twist 分支；每一项只执行一次，并在写回前固定精度与 scale。
  \end{block}
  \source{local 补项：\codepath{lib/coarse_op.cuh:1383-1441}；kernel：\codepath{include/kernels/coarse_op_kernel.cuh:1768-1950}。}
\end{frame}

\begin{frame}[t]{粗场存储协议：(Y) 有方向和 ghost，(X) 是每点 local block}
  \begin{columns}[T,totalwidth=\textwidth]
    \begin{column}{0.54\textwidth}
      \centering
      \begin{tikzpicture}[node distance=0.18cm and 0.2cm]
        \node[flowwide,text width=1.65cm,font=\scriptsize] (q) {粗点 \(q\)};
        \node[state,text width=1.65cm,font=\scriptsize,right=of q] (x) {\(X(q)\)\\\(N_s^cN_c^c\) 方阵};
        \node[state,text width=1.65cm,font=\scriptsize,above right=0.3cm and 0.22cm of q] (yf) {\(Y_{+\mu}(q)\)\\\(q\to q+\mu\)};
        \node[state,text width=1.65cm,font=\scriptsize,below right=0.3cm and 0.22cm of q] (yb) {\(Y_{-\mu}(q)\)\\\(q\to q-\mu\)};
        \node[phase,text width=1.85cm,font=\scriptsize,right=1.45cm of yf] (gf) {forward ghost\\边界邻点};
        \node[phase,text width=1.85cm,font=\scriptsize,right=1.45cm of yb] (gb) {backward ghost\\边界邻点};
        \draw[arrow] (q)--(x);
        \draw[arrow] (q)--(yf);
        \draw[arrow] (q)--(yb);
        \draw[relation] (yf)--(gf);
        \draw[relation] (yb)--(gb);
      \end{tikzpicture}
    \end{column}
    \begin{column}{0.40\textwidth}
      \begin{block}{代码创建的元数据}
        粗 (Y)：full site subset、coarse geometry、nFace=1、双向 link ghost；粗 (X)：scalar geometry、nFace=0、无 ghost。GPU 可能另存 MILC/AoS copy 供 MMA。
      \end{block}
      \begin{block}{方向索引}
        在粗 dslash 中，forward gather 使用 \(Y[d+4](q)\)，backward gather 使用 \(Y[d](q-\hat d)^\dagger\)。
      \end{block}
    \end{column}
  \end{columns}
  \source{字段创建：\codepath{lib/dirac_coarse.cpp:121-170}；方向 gather：\codepath{include/kernels/dslash_coarse.cuh:150-249}。}
\end{frame}

\begin{frame}[fragile,t]{PC 粗化（1/2）：先求 \(X^{-1}\)，再构造 \(\widehat Y\) 的左右乘法}
  \begin{equation*}
    X^{-1}(q)X(q)=I,
    \qquad
    \begin{aligned}
      Yhat_{+\mu}(q)&=X^{-1}(q)Y_{+\mu}(q),\\
      Yhat_{-\mu}(q-\hat\mu)&=Y_{-\mu}(q-\hat\mu)X^{-\dagger}(q).
    \end{aligned}
  \end{equation*}
  \begin{lstlisting}[firstnumber=1]
build_pc_links(Y, X):
    Xinv[q] = batch_inverse(X[q])       # local, independent for each q
    exchange_ghost(Y, both_directions)
    for direction mu and coarse site q:
        Yhat_fwd[q,mu] = Xinv[q] * Y_fwd[q,mu]
        # backward link is stored at q-mu; its right factor is Xinv(q)^dagger
        Yhat_back[q-mu,mu] = Y_back[q-mu,mu] * dagger(Xinv[q])
    exchange_ghost(Yhat, both_directions)
    return Xinv, Yhat
  \end{lstlisting}
  \vspace{0.07cm}
  \begin{alertblock}{最容易写反的地方}
    forward 是左乘 \(X^{-1}\)，backward 是右乘 \(X^{-\dagger}\)，且 backward link 的坐标是邻点 \(q-\mu\)。逐点写成统一的 \(X^{-1}Y\) 会破坏 \(\hat D-A^{-1}D\) 检查。
  \end{alertblock}
  \source{\texttt{\detokenize{calculateYhat}} 的注释与实现：\codepath{lib/coarse_op_preconditioned.in.cu:147-203}；矩阵 tile：\codepath{include/kernels/coarse_op_preconditioned.cuh:48-124}。伪代码保持左右顺序。}
\end{frame}

\begin{frame}[fragile,t]{PC 粗化（2/2）：Schur 补与 QUDA 的 \texttt{\detokenize{DiracCoarsePC::M}} 对齐}
  \begin{columns}[T,totalwidth=\textwidth]
    \begin{column}{0.45\textwidth}
      \begin{equation*}
        D=\begin{pmatrix}A_{ee}&D_{eo}\\D_{oe}&A_{oo}\end{pmatrix},
        \qquad
        S_e=A_{ee}-D_{eo}A_{oo}^{-1}D_{oe}.
      \end{equation*}
      \begin{block}{对称 PC}
        输入只保留目标 parity；\texttt{\detokenize{Dslash}} 负责 \(A^{-1}D\) 的 neighbor hop，\texttt{\detokenize{DslashXpay(...,-1)}} 把两个 hopping 串联并减去单位项，得到目标 parity 的 Schur 动作。
      \end{block}
    \end{column}
    \begin{column}{0.51\textwidth}
      \begin{lstlisting}[firstnumber=1]
apply_schur(in_target, target_parity):
    other = 1 - target_parity
    tmp_other = A_other^{-1} * D_other,target * in_target
    tmp_target = A_target^{-1} * D_target,other * tmp_other
    out = tmp_target - in_target       # QUDA DslashXpay k=-1
    return out

prepare_full_rhs(b_even, b_odd):
    src = A_other^{-1} * b_other
    src = A_target^{-1} * b_target -
          A_target^{-1} * D_target,other * src
    return src

reconstruct_full(x_target, b_other):
    x_other = A_other^{-1} * b_other -
              A_other^{-1} * D_other,target * x_target
    return (x_target, x_other)
      \end{lstlisting}
    \end{column}
  \end{columns}
  \source{\texttt{\detokenize{DiracCoarsePC::M}}：\codepath{lib/dirac_coarse.cpp:530-560}；prepare/reconstruct：\codepath{lib/dirac_coarse.cpp:570-626}。公式是相同 block elimination 的数学写法。}
\end{frame}

\begin{frame}[fragile,t]{粗 dslash 的单点动作：先 forward/backward gather，再统一乘 \((-\kappa)\)}
  \begin{equation*}
    (D_{\rm hop}\psi)(q)=
    \sum_{\mu=0}^{3}\left[
      Y_{+\mu}(q)\psi(q+\hat\mu)+
      Y_{-\mu}^{\dagger}(q-\hat\mu)\psi(q-\hat\mu)
    \right],
  \end{equation*}
  \begin{lstlisting}[firstnumber=1]
coarse_dslash(out, in, parity, use_X):
    for q in destination_sites(parity):
        acc = zero_vector(q)
        for mu in 0..3:
            qf = neighbor(q, +mu)
            qb = neighbor(q, -mu)
            psi_f = halo_or_local(in, qf, parity^1)
            psi_b = halo_or_local(in, qb, parity^1)
            acc += Y_fwd[q,mu] * psi_f
            acc += dagger(Y_back[qb,mu]) * psi_b
        acc *= -kappa
        if use_X:
            acc += X[q] * in[q, parity]
        out[q, parity] = acc
  \end{lstlisting}
  \vspace{0.05cm}
  \begin{block}{实现对齐点}
    内部点直接从 \texttt{\detokenize{inA}} 读；通信边界从 ghost 读。粗 \(X\) 的作用与 hopping 在同一 kernel 的最终阶段相加；\texttt{\detokenize{Dslash}} 只保留 hopping，\texttt{\detokenize{DslashXpay/M}} 再组合 local 项。
  \end{block}
  \source{kernel 的 forward/backward gather：\codepath{include/kernels/dslash_coarse.cuh:126-255}；local \(X\)：\codepath{include/kernels/dslash_coarse.cuh:257-290}；\(-\kappa\) 汇总：\codepath{include/kernels/dslash_coarse.cuh:292-324}。}
\end{frame}

\begin{frame}[fragile,t]{MPI/halo 复现（1/2）：通信只改变 neighbor 数据的取得方式}
  \begin{columns}[T,totalwidth=\textwidth]
    \begin{column}{0.49\textwidth}
      \begin{lstlisting}[firstnumber=1]
exchange_spinor_halo(in, depth=1):
    for mu in dimensions:
        pack plus_face  = in[coord[mu] == L_local[mu]-1]
        pack minus_face = in[coord[mu] == 0]
        include parity, spin, color, rhs_batch
        post Irecv(plus_ghost  <- rank_minus_mu)
        post Irecv(minus_ghost <- rank_plus_mu)
        post Isend(plus_face  -> rank_plus_mu)
        post Isend(minus_face -> rank_minus_mu)
    return requests, ghost_buffers

apply_coarse_dslash:
    requests, ghost = exchange_spinor_halo(in)
    compute_interior_sites(out, in)       # overlap
    wait_all(requests)
    compute_boundary_sites(out, in, ghost)
      \end{lstlisting}
    \end{column}
    \begin{column}{0.46\textwidth}
      \begin{block}{每个粗 dslash 只需一层}
        Wilson/Clover coarse hopping 是一跳；因此 \(nFace=1\)。ASQTAD 细格 setup 的 UV/LV 需要 \texttt{\detokenize{nFace=3}}，不要把这个 setup halo 深度误用于粗 dslash。
      \end{block}
      \begin{block}{单 rank 极限}
        \texttt{\detokenize{rank_plus_mu=rank_minus_mu=self}} 时用周期坐标直接索引即可；这应与 MPI 版本共享同一 \texttt{\detokenize{neighbor()}} 逻辑。
      \end{block}
    \end{column}
  \end{columns}
  \source{粗 dslash 的 halo/bulk 分支：\codepath{include/kernels/dslash_coarse.cuh:148-255}；coarse field nFace/ghost：\codepath{lib/dirac_coarse.cpp:130-148}；setup nFace：\codepath{lib/coarse_op.cuh:1040-1057}。伪代码给出 MPI 后端契约。}
\end{frame}

\begin{frame}[t]{MPI/halo 复现（2/2）：链接 halo、旋量 halo 和聚合边界要分开}
  \begin{center}
    \begin{tikzpicture}[node distance=0.25cm and 0.24cm]
      \node[flowwide,text width=2.0cm,font=\scriptsize] (v) {setup 前\\交换 (V)};
      \node[flowwide,text width=2.0cm,font=\scriptsize,right=of v] (uv) {UV/LV\\读取邻点 (V)};
      \node[flowwide,text width=2.0cm,font=\scriptsize,right=of uv] (y) {写 (Y/X)\\并交换 link ghost};
      \node[flowwide,text width=2.0cm,font=\scriptsize,below=0.38cm of y] (s) {solve 前\\交换旋量};
      \node[flowwide,text width=2.0cm,font=\scriptsize,left=of s] (d) {粗 dslash\\boundary gather};
      \node[flowwide,text width=2.0cm,font=\scriptsize,left=of d] (r) {R/P\\仅按 map 聚合};
      \draw[arrow] (v)--(uv)--(y);
      \draw[arrow] (y)--(s)--(d)--(r);
      \draw[relation] (r.north) |- node[left,font=\scriptsize]{下一层 coarse field} (v.south);
    \end{tikzpicture}
  \end{center}
  \vspace{0.15cm}
  \begin{columns}[T,totalwidth=\textwidth]
    \begin{column}{0.48\textwidth}
      \begin{block}{setup 期间}
        \texttt{\detokenize{calculateY}} 先让 \(V\) 交换 ghost；跨 rank 的 \(V^\dagger H V\) 才能得到正确边界 coarse link。写好 \(Y\) 后，PC \texttt{\detokenize{calculateYhat}} 还要让 \(Y\) 和 \(\widehat Y\) 的双向 link ghost 有效。
      \end{block}
    \end{column}
    \begin{column}{0.48\textwidth}
      \begin{alertblock}{aggregate 不跨 rank 的假设不能擅自删除}
        map 是 local volume 上的聚合索引；若 block 跨 MPI 分区，需要先重新定义 aggregate ownership 或通信，不能仅靠 coarse dslash 的 halo 修补 setup 的 \(V\) 访问。
      \end{alertblock}
    \end{column}
  \end{columns}
  \source{V ghost：\codepath{lib/coarse_op.cuh:1056-1062}；Y link ghost：\codepath{lib/coarse_op_preconditioned.in.cu:197-203}、\codepath{lib/dirac_coarse.cpp:260-270}；粗 dslash boundary：\codepath{include/kernels/dslash_coarse.cuh:156-178, 211-231}。}
\end{frame}

\begin{frame}[fragile,t]{递归 V-cycle（1/2）：按 QUDA 的 residual 分支逐行展开}
  \begin{lstlisting}[firstnumber=1]
MG(level, x, b):
    parity = target_parity_from_matpc_type()
    outer = MAT_SOLUTION if b.is_full else MATPC_SOLUTION
    inner = level.coarse_grid_solution_type
    uniform = (level.smoother_solve_type == DIRECT_PC and inner == MATPC)
           or (level.smoother_solve_type == DIRECT    and inner == MAT)
    assert not (outer == MATPC and inner == MAT)
    assert not (inner == MATPC and smoother_type != DIRECT_PC)

    if level is coarsest:
        (out, in) = dirac.prepare(x, b, outer)
        if presmoother exists: presmoother(out, in)
        dirac.reconstruct(x, b, outer)
        return x

    (out, in) = dirac.prepare(x, b, outer)
    if presmoother exists: presmoother(out, in)
    if not uniform: dirac.reconstruct(x, b, inner)

    if presmoother exists and uniform:
        residual = presmoother.returned_residual()
    else if no presmoother and uniform:
        residual = in
    else if no presmoother:
        residual = b
    else:
        r = b - A_level*x
        residual = r
  \end{lstlisting}
  \vspace{0.04cm}
  \begin{block}{到这里为止的含义}
    预平滑可能返回已经计算好的 (b-Ax)；只有 smoother solution type 与 coarse correction type 不一致时，代码才重构 (x) 后显式计算真 residual。
  \end{block}
  \source{真实控制流：\codepath{lib/multigrid.cpp:1131-1183}；\texttt{\detokenize{uniform}} 判定与不支持组合：同文件 1146-1169。伪代码按分支展开。}
\end{frame}

\begin{frame}[fragile,t]{递归 V-cycle（2/2）：限制、粗解、延拓、后平滑和最终重构}
  \begin{lstlisting}[firstnumber=1]
    if transfer exists:
        r_c = R(residual)
        x_c = allocate_zero_coarse_like(r_c)
        coarse_solver(x_c, r_c)             # recurse or bottom Krylov

        if no presmoother:
            if inner == outer:
                P_add(x, x_c)
            else:
                P_add(x[target_parity], x_c)
        else:
            correction = zero_like(r)
            if inner == outer:
                P(correction, x_c)
                x += correction
            else:
                P(correction[target_parity], x_c)
                x[target_parity] += correction[target_parity]

    if not uniform:
        (out, in) = dirac.prepare(x, b, inner)
    if postsmoother exists:
        postsmoother(out, in)                # use_init_guess = YES
    dirac.reconstruct(x, b, outer)
    return x
  \end{lstlisting}
  \vspace{0.06cm}
  \begin{columns}[T,totalwidth=\textwidth]
    \begin{column}{0.48\textwidth}
      \begin{block}{为什么要 \texttt{\detokenize{P_add}}}
        有 pre-smoother 时，粗解是对当前平滑后 \(x\) 的 correction；必须 \(x\leftarrow x+P x_c\)，不能覆盖掉 \(x\)。
      \end{block}
    \end{column}
    \begin{column}{0.48\textwidth}
      \begin{alertblock}{粗层没有 transfer}
        代码仍执行本层 prepare、可选 pre-smooth、reconstruct；只在 `transfer` 存在时进入 \(R\to\) coarse solver \(\to P\)。
      \end{alertblock}
    \end{column}
  \end{columns}
  \source{限制/递归/延拓/post：\codepath{lib/multigrid.cpp:1184-1221}；pre/post solver 的初值与 residual 标志：\codepath{lib/multigrid.cpp:273-337}。}
\end{frame}

\begin{frame}[t]{V-cycle 的语义分层：平滑器、coarse solver 和 cycle wrapper 不是一回事}
  \begin{columns}[T,totalwidth=\textwidth]
    \begin{column}{0.31\textwidth}
      \begin{block}{平滑器}
        \(S_{pre}\) 和 \(S_{post}\) 是本层局部迭代；\texttt{\detokenize{nu_pre}} 影响粗化前的误差，\texttt{\detokenize{nu_post}} 清理粗校正带来的高频误差。
      \end{block}
    \end{column}
    \begin{column}{0.31\textwidth}
      \begin{block}{coarse solver}
        \texttt{\detokenize{coarse_solver}} 接收 \(r_c\)，输出 \(x_c\)。非底层可以直接调用下层 MG；底层可以是 GCR/CG 等明确配置的 solver。
      \end{block}
    \end{column}
    \begin{column}{0.31\textwidth}
      \begin{block}{cycle type}
        \texttt{\detokenize{VCYCLE}} 在允许的层级直接把下层 MG 作为 coarse solver；\texttt{\detokenize{RECURSIVE}} 创建 \texttt{\detokenize{PreconditionedSolver}} wrapper，把下层 MG 放进 Krylov 预条件链。
      \end{block}
    \end{column}
  \end{columns}
  \vspace{0.16cm}
  \begin{center}
    \begin{tikzpicture}[node distance=0.20cm and 0.22cm]
      \node[flowwide] (a) {本层 \(r\)};
      \node[flowwide,right=of a] (b) {pre-smooth};
      \node[flowwide,right=of b] (c) {R + 下层 MG};
      \node[flowwide,right=of c] (d) {P + add};
      \node[flowwide,right=of d] (e) {post-smooth};
      \draw[arrow] (a)--(b)--(c)--(d)--(e);
      \node[phase,below=0.32cm of c,text width=2.65cm] (w) {recursive wrapper 可在更外层 GCR 中多次调用该链};
      \draw[relation] (w)--(c);
    \end{tikzpicture}
  \end{center}
  \source{cycle wrapper：\codepath{lib/multigrid.cpp:535-702}；运行期 V-cycle：\codepath{lib/multigrid.cpp:1171-1221}。}
\end{frame}

\begin{frame}[fragile,t]{Schwarz/MR 平滑器：每次 MR 更新都必须保留复数 α 和残差状态}
  \begin{equation*}
    \alpha_k=\frac{(A r_k,r_k)}{(A r_k,A r_k)},
    \qquad
    x_{k+1}=x_k+\omega\alpha_k r_k,
    \qquad
    r_{k+1}=r_k-\omega\alpha_k A r_k.
  \end{equation*}
  \begin{lstlisting}[firstnumber=1]
MR_smooth(x, b, maxiter, omega, Nsteps, schwarz):
    r = b - A*x if use_init_guess else b;  x = 0 if no_guess
    for step in 0 .. Nsteps-1:
        normalize = sqrt(initial_residual_norm2)
        rs = r / normalize
        k = 0
        while k < maxiter and norm2(r) > delta2:
            Ar = A_sloppy(rs)
            alpha = dot(Ar, rs) / dot(Ar, Ar)   # complex
            xs += omega * alpha * rs
            rs -= omega * alpha * Ar
            k += 1
        x += normalize * xs
        if compute_true_res or Nsteps > 1:
            r = b - A*x
        else:
            r = normalize * rs
        if step+1 == Nsteps or norm2(r) < stopping_threshold:
            break
    return x, r
  \end{lstlisting}
  \vspace{0.03cm}
  \begin{block}{QUDA 的额外状态}
    \texttt{\detokenize{return_residual=true}} 只给 pre-smoother 使用；\texttt{\detokenize{sloppy_converge=true}} 允许内部 sloppy residual 停止；block Schwarz 时可改为局部 reduction，并按 red/black active mask 更新。
  \end{block}
  \source{MR 主循环：\codepath{lib/inv_mr_quda.cpp:65-207}；pre-smoother \texttt{\detokenize{return_residual}}：\codepath{lib/multigrid.cpp:279-319}。伪代码保留缩放、真残差和 Nsteps 语义。}
\end{frame}

\begin{frame}[fragile,t]{Schwarz 的作用：改变局部预条件器，不改变 MG 的 (R/P) 逻辑}
  \begin{columns}[T,totalwidth=\textwidth]
    \begin{column}{0.47\textwidth}
      \begin{block}{Additive Schwarz}
        每次 active 块都用局部数据更新，两个颜色可同时准备，块间 hopping 暂时关闭；适合把局部块问题作为 MR 内部预条件。
      \end{block}
      \begin{block}{Multiplicative Schwarz}
        red/black 交替更新，源码要求 \texttt{\detokenize{Nsteps}} 为偶数；全局网格未做 block Schwarz 时还按 node parity 交替。
      \end{block}
    \end{column}
    \begin{column}{0.49\textwidth}
      \begin{lstlisting}[firstnumber=1]
if block_schwarz:
    if type == MULTIPLICATIVE:
        active = RED if step even else BLACK
        set_DD_active(active)
    else:                       # additive
        set_DD_active(RED | BLACK, no_block_hopping)
    disable_global_reduction_if_blocks_local()

run_MR_inner_steps()

if block_schwarz:
    reset_DD_masks()
      \end{lstlisting}
      \begin{alertblock}{setup 例外}
        普通 Krylov solver 不支持该 Schwarz 配置时，QUDA 在 null-vector finding 期间把 \texttt{\detokenize{commDim}} 设为全 1、临时禁用 Schwarz，结束后恢复。
      \end{alertblock}
    \end{column}
  \end{columns}
  \source{MR Schwarz 分支：\codepath{lib/inv_mr_quda.cpp:106-179}；multiplicative 偶数限制：同文件 15-20；setup 临时禁用：\codepath{lib/multigrid.cpp:1321-1329, 1460-1466}。}
\end{frame}

\begin{frame}[fragile,t]{外层 GCR：把一次 MG 应用当作可变预条件器 (K(r))}
  \begin{lstlisting}[firstnumber=1]
flexible_GCR(A, b, K=MG, restart, maxiter):
    x = initial_guess_or_zero(); r = b - A*x
    basis_p = []; basis_Ap = []
    while not converged(r) and iter < maxiter:
        p = K(r)                         # one MG V-cycle / wrapper
        Ap = A(p)
        for j in 0 .. len(basis_Ap)-1:
            beta = dot(basis_Ap[j], Ap) / dot(basis_Ap[j], basis_Ap[j])
            Ap -= beta * basis_Ap[j]
            p  -= beta * basis_p[j]
        alpha = dot(Ap, r) / dot(Ap, Ap)
        x += alpha * p
        r -= alpha * Ap
        append(basis_p, p); append(basis_Ap, Ap)
        if len(basis_p) == restart or reliable_update_needed(r):
            r = b - A*x                    # true residual
            basis_p.clear(); basis_Ap.clear()
    return x
  \end{lstlisting}
  \vspace{0.04cm}
  \begin{columns}[T,totalwidth=\textwidth]
    \begin{column}{0.48\textwidth}
      \begin{block}{为什么是 flexible}
        sloppy 精度、Schwarz、递归粗解和 setup 状态可能让每次 \(K\) 的数值行为不同；GCR/Flexible GCR 通过每轮显式保留 \(p_k\) 和 \(Ap_k\) 处理这种非固定预条件效果。
      \end{block}
    \end{column}
    \begin{column}{0.48\textwidth}
      \begin{alertblock}{可靠更新}
        QUDA 在 Krylov 满、达到阈值或 \(r_k\) 相对下降过快时更新解并重算真 residual；若真残差上升过多会提前告警退出。
      \end{alertblock}
    \end{column}
  \end{columns}
  \source{GCR 预条件调用与正交化：\codepath{lib/inv_gcr_quda.cpp:303-358}；可靠更新/真残差：同文件 348-407；MG 自举把自身交给 GCR：\codepath{lib/multigrid.cpp:1340-1363}。伪代码是等价 Krylov 逻辑。}
\end{frame}

\begin{frame}[t]{精度、MMA 和存储布局：只替换“如何算”，不替换“算什么”}
  \begin{columns}[T,totalwidth=\textwidth]
    \begin{column}{0.24\textwidth}
      \begin{block}{precision\_null}
        决定 (B/V) 及粗 link 的存储精度；CPU setup 小于单精度时源码提升到至少 single。
      \end{block}
    \end{column}
    \begin{column}{0.24\textwidth}
      \begin{block}{two-pass}
        低精度 (V) 先求幅值、设 scale，再重算块正交；scale 必须随字段元数据传播。
      \end{block}
    \end{column}
    \begin{column}{0.24\textwidth}
      \begin{block}{MMA}
        setup/transfer/dslash 各自有开关；MMA 只改变 tile、order 和 kernel 调度，矩阵乘的共轭/转置语义不能变。
      \end{block}
    \end{column}
    \begin{column}{0.24\textwidth}
      \begin{block}{halo precision}
        smoother halo 可独立于 (V/Y) 精度；setup 发现 quarter halo 时暂时提升到 half。
      \end{block}
    \end{column}
  \end{columns}
  \vspace{0.22cm}
  \begin{center}
    \begin{tikzpicture}[node distance=0.2cm and 0.27cm]
      \node[flowwide] (math) {复数矩阵语义\(RDP\)};
      \node[flowwide,right=of math] (layout) {field order\\AoS/SoA/tile};
      \node[flowwide,right=of layout] (prec) {precision\\scale/convert};
      \node[flowwide,right=of prec] (comm) {parallel\\halo/reduction};
      \draw[arrow] (math)--(layout)--(prec)--(comm);
    \end{tikzpicture}
  \end{center}
  \begin{alertblock}{迁移顺序}
    先用 double/single CPU reference 固定数学结果，再逐项打开 parity、MPI、half/fixed-point、MMA；每一步均回到 \(RP\)、\(D_c-RDP\) 和 PC dslash 三个不变量。
  \end{alertblock}
    \source{精度与 setup 字段：\codepath{include/quda.h:647-669, 755-781}；quarter halo 临时提升：\codepath{lib/multigrid.cpp:1331-1333}；MMA/字段 copy：\codepath{lib/dirac_coarse.cpp:243-296}。}
\end{frame}

\begin{frame}[t]{参数到行为映射（1/2）：几何、零空间和 setup 参数}
  \scriptsize
  \begin{tabularx}{\textwidth}{@{}l l X@{}}
    \toprule
    参数 & 所在阶段 & 实现行为与复现要求 \\
    \midrule
    \texttt{\detokenize{n_level}} & 层级 & 创建 \(MG_0,\ldots,MG_{n-1}\)；不得超过编译时最大递归深度。 \\
    \texttt{\detokenize{n_vec}} & (B/V) & 每层粗颜色列数；必须不超过 aggregate 容量，且决定粗算子矩阵大小。 \\
    \texttt{\detokenize{n_vec_batch}} & null setup & 每批同时生成的 \(B_i\) 数；\texttt{len(B) \% batch == 0}，不整除直接失败。 \\
    \texttt{\detokenize{geo_block_size}} & map/setup & 逐维 \(b_\mu\)；影响粗体积、局部容量、long-link 是否可安全粗化。 \\
    \texttt{\detokenize{spin_block_size}} & map/setup & \(s_f\to s_c\) 的整数分块；staggered 用 0 并映射成 2 个 coarse spin。 \\
    \texttt{\detokenize{n_block_ortho}} & V & 每个 aggregate 重复块 Gram--Schmidt 的次数。 \\
    \texttt{\detokenize{block_ortho_two_pass}} & V 精度 & half/fixed-point 下先量幅值再重新计算正交化。 \\
    \texttt{\detokenize{precision_null}} & B/V/Y & null vector 和 coarse link 的存储/计算精度来源；CPU 小精度有提升规则。 \\
    \texttt{\detokenize{setup_inv_type}} & B & CG/CA-CG 走 \(D^\dagger D\)，MG 走 GCR+MG 自举，其他按 solver 工厂创建。 \\
    \texttt{\detokenize{setup_maxiter}}, \texttt{\detokenize{setup_tol}} & B & 每次 setup solve 的最大迭代和相对残差目标；refresh 使用独立 maxiter。 \\
    \texttt{\detokenize{num_setup_iter}} & B & 每层重复“正交化→分批 solve→正交化→重建/传播”的次数。 \\
    \bottomrule
  \end{tabularx}
  \vspace{0.08cm}
  \begin{block}{配置解读顺序}
    先检查 \texttt{\detokenize{geo_block_size/spin_block_size/n_vec}} 的维数合法性，再决定 setup solver 和精度；顺序反过来会得到可运行但无法形成有效粗空间的配置。
  \end{block}
  \source{参数声明：\codepath{include/quda.h:638-701}；每层拷贝：\codepath{include/multigrid.h:181-217}；setup 使用：\codepath{lib/multigrid.cpp:1279-1319, 1383-1413}。}
\end{frame}

\begin{frame}[t]{参数到行为映射（2/2）：运行期、PC、cycle 和验证}
  \scriptsize
  \begin{tabularx}{\textwidth}{@{}l l X@{}}
    \toprule
    参数 & 所在阶段 & 实现行为与复现要求 \\
    \midrule
    \texttt{\detokenize{nu_pre}} / \texttt{\detokenize{nu_post}} & V-cycle & 分别创建 pre/post solver；pre 返回 residual，post 使用当前 \(x\) 作为初值。 \\
    \texttt{\detokenize{smoother}} & smooth & MR/GCR/CA 等实际平滑器；MR 的 \(\omega\) 更新与 Schwarz 需单独实现。 \\
    \texttt{\detokenize{smoother_solve_type}} & prepare/粗层 & \texttt{\detokenize{DIRECT}} 或 \texttt{\detokenize{DIRECT_PC}}；PC coarse solution 必须配 \texttt{\detokenize{DIRECT_PC}} smoother。 \\
    \texttt{\detokenize{coarse_grid_solution_type}} & R/P 接口 & \texttt{\detokenize{MAT}} 传 full field；\texttt{\detokenize{MATPC}} 只传目标 parity，并决定是否调用 prepare/reconstruct。 \\
    \texttt{\detokenize{matpc_type}} & parity/Schur & 由 \texttt{\detokenize{EVEN_EVEN}} / \texttt{\detokenize{ODD_ODD}} 选择目标 parity；PC 的 \(S_e/S_o\) 作用于 single-parity field。 \\
    \texttt{\detokenize{cycle_type}} & coarse solver & \texttt{\detokenize{VCYCLE}} 可直接链接下层 MG；\texttt{\detokenize{RECURSIVE}} 通过 \texttt{\detokenize{PreconditionedSolver}} 接到 Krylov。 \\
    \texttt{\detokenize{use_init_guess}} & 各 solver & null setup=YES；pre smoother=NO；post smoother=YES；coarse solver 默认 NO，deflation+\texttt{\detokenize{coarse_guess}} 可打开。 \\
    \texttt{\detokenize{run_verify}} & setup 后 & 检查 \(PP^\dagger B\)、\(P^\dagger P\)、\(D_c-RDP\)、PC dslash 和 normality。 \\
    \texttt{\detokenize{generate_all_levels}} & setup 传播 & YES：下一层独立生成；NO：\(B_{\ell+1}=R B_\ell\) 后重建 Transfer。 \\
    \texttt{\detokenize{coarse_solver_*}} & 粗解 & 给递归 wrapper 的 inverter、tol、maxiter、CA basis；底层显式 Krylov 仍有自己的停止条件。 \\
    \texttt{\detokenize{smoother_schwarz_*}} & smooth & additive/multiplicative、cycle 数、halo/reduction 行为；setup 普通 Krylov 期间可能被临时关闭。 \\
    \bottomrule
  \end{tabularx}
  \vspace{0.06cm}
  \begin{alertblock}{关键组合限制}
    公共入口当前只接受 direct outer solve；MG smoother solve type 只支持 direct/direct-PC；\texttt{\detokenize{outer=MATPC}} 且 \texttt{\detokenize{inner=MAT}} 是源码显式拒绝的组合。
  \end{alertblock}
  \source{参数声明：\codepath{include/quda.h:731-835}；smoother 创建：\codepath{lib/multigrid.cpp:273-337}；组合检查：\codepath{lib/multigrid.cpp:1146-1169}、\codepath{lib/interface_quda.cpp:2863-2871}。}
\end{frame}

\begin{frame}[fragile,t]{从零实现的端到端伪代码：先单 rank，再逐层打开复杂度}
  \begin{lstlisting}[firstnumber=1]
main(cfg):
    U, C = load_gauge_and_optional_clover(cfg)
    D = make_fine_dirac(U, C, cfg.mass, cfg.kappa, cfg.mu)
    levels = []
    B = initialize_null_vectors(cfg.level[0])

    for level in 0 .. cfg.n_level-1:
        L = setup_level(level, D, B, cfg.level[level])
        levels.append(L)
        if level + 1 < cfg.n_level:
            if cfg.generate_all_levels:
                B = initialize_null_vectors(cfg.level[level+1])
            else:
                B = [L.R(v) for v in L.B]
            D = L.coarse_operator()             # Y/X or Yhat/Xinv

    assert verify_all_levels(levels, D, cfg)
    x = zero_like(rhs)
    for outer_iter in 1 .. cfg.outer_maxiter:
        correction = MG(levels[0], zero_like(rhs), rhs)
        update_outer_solution(x, correction, cfg.outer_solver)
        if true_norm(rhs - D_fine(x)) <= cfg.outer_tol:
            break
    return x
  \end{lstlisting}
  \vspace{0.05cm}
  \begin{columns}[T,totalwidth=\textwidth]
    \begin{column}{0.48\textwidth}
      \begin{block}{更精确的外层}
        实际 QUDA 通常把 MG 作为 GCR/其他 solver 的 preconditioner；上面的 \texttt{\detokenize{update_outer_solution}} 应替换为上一页的 flexible GCR，而不是简单累加多个 V-cycle。
      \end{block}
    \end{column}
    \begin{column}{0.48\textwidth}
      \begin{alertblock}{粗算子持久化}
        若要跨运行保存 null vectors/coarse operator，必须同时保存几何、spin map、精度/scale、Y/X 方向约定和 gauge 边界条件，否则数组本身不足以复现动作。
      \end{alertblock}
    \end{column}
  \end{columns}
  \source{公共 API 生命周期：\codepath{include/quda.h:1256-1289}；测试中的构造/solve/update：\codepath{tests/multigrid_evolve_test.cpp:308-365}。伪代码是脱离 QUDA 的组合实现。}
\end{frame}

\begin{frame}[t]{验证矩阵（1/2）：把“构造正确”拆成五个可定位的误差源}
  \scriptsize
  \begin{tabularx}{\textwidth}{@{}l l X X@{}}
    \toprule
    编号 & 随机对象 & 计算 & 失败时优先定位 \\
    \midrule
    V1 & 细 \(B_i\) & \(e_i=\|B_i-P R B_i\|/\|B_i\|\) & map、spin map、块正交化、V layout \\
    V2 & 粗 \(\eta\) & \(e_c=\|\eta_c-RP\eta_c\|/\|\eta_c\|\) & restrict 清零/累加、coarse parity、V 共轭 \\
    V3 & 粗随机 \(\eta\) & \(e_D=\|(D_c-RDP)\eta\|/\|RDP\eta\|\) & Y/X 分类、\(\gamma\) 投影、\(\kappa\)、clover/mass \\
    V4 & 单 parity \(\phi\) & \(e_{PC}=\|\hat D-A^{-1}D\|\phi/\|\cdot\|\) & Xinv 逆、Yhat 左右乘、ghost 方向 \\
    V5 & 任意 \(\eta\) & \(\operatorname{Im}(\eta^\dagger M^\dagger M\eta)/\operatorname{Re}(\cdot)\) & dagger、spin/color 转置、parity 选择 \\
    \bottomrule
  \end{tabularx}
  \vspace{0.14cm}
  \begin{columns}[T,totalwidth=\textwidth]
    \begin{column}{0.49\textwidth}
      \begin{block}{精度相关阈值（源码）}
        quarter/half：\(5\times10^{-2}\)；single：\(2\times10^{-3}\)；更高精度默认 \(10^{-8}\)。这是 QUDA verify 的检查阈值，不是所有物理配置的收敛保证。
      \end{block}
    \end{column}
    \begin{column}{0.47\textwidth}
      \begin{alertblock}{特殊跳过}
        ASQTAD 长链接或过小 aggregate 会令 native coarse 与显式投影不完全对应；源码按 transfer type、operator 和 block size 条件跳过部分 verify，报告时必须保留该状态。
      \end{alertblock}
    \end{column}
  \end{columns}
  \source{阈值与 V1/V2：\codepath{lib/multigrid.cpp:749-880}；V3/V4/V5 与 ASQTAD skip：\codepath{lib/multigrid.cpp:883-1128}。}
\end{frame}

\begin{frame}[fragile,t]{验证矩阵（2/2）：每个测试都要记录范数、相对误差和配置}
  \begin{lstlisting}[firstnumber=1]
verify_level(L):
    for i, B_i in enumerate(L.B):
        coarse = L.R(B_i); projected = L.P(coarse)
        report("null_span", i, norm2(B_i-projected)/norm2(B_i))

    eta = random_coarse_field(L)
    report("coarse_span", norm2(eta-L.R(L.P(eta)))/norm2(eta))

    eta = random_coarse_field(L)
    explicit = L.R(fine_D(L.P(eta)))
    native = L.coarse_apply(eta)
    report("galerkin", norm2(native-explicit)/norm2(explicit))

    if L.has_pc:
        phi = random_single_parity_field(L)
        report("pc_dslash", norm2(L.pc_dslash(phi)-
              L.Xinv * L.unpreconditioned_dslash(phi)))

    assert no_nan_or_inf(all_fields_and_scalars)
  \end{lstlisting}
  \vspace{0.07cm}
  \begin{block}{记录格式}
    每行至少保存：level、fine/coarse lattice、\(b_\mu\)、\(Nvec\)、spin block、precision、parity、operator type、范数分母、L2 relative、max deviation。只报告“通过”而没有这些上下文，无法复现误差。
  \end{block}
  \source{QUDA verify 的随机向量与对照动作：\codepath{lib/multigrid.cpp:785-832, 866-928, 930-1068}；伪代码为测试 harness 抽象。}
\end{frame}

\begin{frame}[t]{端到端复现阶梯：每一级都比上一级多一个可控变量}
  \begin{center}
    \begin{tikzpicture}[node distance=0.22cm and 0.18cm]
      \node[phase,text width=1.75cm] (s1) {1. CPU 单 rank\\double/single\\显式 (RDP)};
      \node[phase,text width=1.75cm,right=of s1] (s2) {2. CPU 单 rank\\QUDA 风格\\Y/X + MR};
      \node[phase,text width=1.75cm,right=of s2] (s3) {3. 单 GPU\\checkerboard\\低精度};
      \node[phase,text width=1.75cm,right=of s3] (s4) {4. MPI\\spinor/link\\halo};
      \node[phase,text width=1.75cm,right=of s4] (s5) {5. 多层\\PC/recursive\\外层 GCR};
      \draw[arrow] (s1)--(s2)--(s3)--(s4)--(s5);
    \end{tikzpicture}
  \end{center}
  \vspace{0.18cm}
  \scriptsize
  \begin{tabularx}{\textwidth}{@{}l X X@{}}
    \toprule
    阶段 & 必须通过 & 不能提前声称 \\
    \midrule
    单 rank reference & V1--V3、显式 coarse apply 与 \texttt{\detokenize{RDP}} 一致 & GPU/MPI 性能 \\
    单 GPU & parity、field order、precision scale、V1--V5 & MPI overlap \\
    MPI & 单 rank 与多 rank结果、ghost 边界、全局 reduction & 任意 aggregate 跨 rank \\
    多层/PC & 每层 V1--V5、Schur prepare/reconstruct、外层真残差 & 只看 sloppy residual \\
    参数扫描 & 固定 RHS/初值/硬件后比较 \(Nvec,b,\nu\) & 无基线的“加速” \\
    \bottomrule
  \end{tabularx}
  \vspace{0.12cm}
  \begin{alertblock}{停止条件}
    在当前层的所有不变量未通过前，不进入下一层。这样失败时能区分“数学构造错误”“通信错误”“精度错误”和“求解器参数错误”。
  \end{alertblock}
  \source{动态测试入口可参考：\codepath{tests/multigrid_evolve_test.cpp:103-179, 308-365}、\codepath{tests/multigrid_benchmark_test.cpp:135-153}；本页为复现计划，非本次实测数据。}
\end{frame}

\begin{frame}[t]{失败定位表：从哪一个不变量失败，就回到哪一段伪代码}
  \scriptsize
  \begin{tabularx}{\textwidth}{@{}l l X@{}}
    \toprule
    症状 & 根因候选 & 最小定位实验 \\
    \midrule
    \(RP\eta\) 偏离大 & map 顺序、parity 或 restrict 未清零 & 关闭 GPU/MPI，逐 aggregate 打印 \(V^\dagger V-I\) \\
    \(B-PRB\) 偏离大 & 块正交化列顺序、spin map、two-pass scale & \texttt{\detokenize{n_block_ortho=1}}、double reference，对比每列归一范数 \\
    \(D_c-RDP\) 偏离大 & internal/external 分类、\(\gamma\) 符号、\(\kappa\) 重复/遗漏 & 用单一 point coarse basis，逐方向比较 \(Y/X\) \\
    PC dslash 偏离大 & \(X^{-1}\) 左右乘反、backward 坐标错 & 只测一个 parity、一个方向，打印 \(Yhat_{\rm fwd/back}\) \\
    单 rank 对但 MPI 错 & ghost face index、parity pack、neighbor rank & 只激活一个分区方向，比较边界点逐项结果 \\
    残差下降后上升 & sloppy residual 当真 residual、GCR 更新/重构错 & 每次 cycle 计算 (b-Ax)，禁用可靠更新优化 \\
    setup 运行但粗解无效 & (B) 未按配置传播、nvec/batch 不匹配 & 保存每层 (B,V)，逐层执行 V1--V3 \\
    \bottomrule
  \end{tabularx}
  \vspace{0.18cm}
  \begin{center}
    \begin{tikzpicture}[node distance=0.18cm and 0.22cm]
      \node[flowwide] (a) {先测 (RP)};
      \node[flowwide,right=of a] (b) {再测 (RDP)};
      \node[flowwide,right=of b] (c) {再测 PC/halo};
      \node[flowwide,right=of c] (d) {最后看 solver residual};
      \draw[arrow] (a)--(b)--(c)--(d);
    \end{tikzpicture}
  \end{center}
  \source{失败模式对应源码：\codepath{lib/multigrid.cpp:745-1128}、\codepath{lib/multigrid.cpp:1131-1224}、\codepath{include/kernels/dslash_coarse.cuh:126-324}；表为调试决策抽象。}
\end{frame}

\begin{frame}[t]{源码控制流对照（1/2）：setup 与层级装配的逐段映射}
  \scriptsize
  \begin{tabularx}{\textwidth}{@{}l l X@{}}
    \toprule
    复现伪代码阶段 & QUDA 实现 & 读者应提取的规则 \\
    \midrule
    初始化/选择 (B) & \codepath{lib/multigrid.cpp:36-62} & random、file、eig、iterative、free-field 按条件分支；不要混淆 \texttt{\detokenize{compute_null_vector}} 与 \texttt{\detokenize{vec_load}}。 \\
    创建 level & \codepath{lib/multigrid.cpp:72-186} & transfer/coarse fields/next MG/coarse solver 有严格依赖顺序。 \\
    创建 smoother & \codepath{lib/multigrid.cpp:273-337} & pre 返回 residual、post 使用 init guess；精度、Schwarz、tol 在此装配。 \\
    创建 coarse Dirac & \codepath{lib/multigrid.cpp:342-434} & 根据 \texttt{\detokenize{coarse_grid_solution_type}} 与 smoother solve type 选择 coarse 或 coarse-PC。 \\
    setup solver & \codepath{lib/multigrid.cpp:1275-1363} & maxiter/tol/use-init/batch/CG vs GCR vs MG 自举。 \\
    batch + 正交 & \codepath{lib/multigrid.cpp:1365-1427} & pre/post global orthonormalize、\texttt{\detokenize{TEST_VECTOR}} 分支、每批写回 (B)。 \\
    向下传播 & \codepath{lib/multigrid.cpp:1429-1451} & generate-all 或 (R B)，随后 reset 下一层。 \\
    \bottomrule
  \end{tabularx}
  \vspace{0.14cm}
  \begin{block}{如何使用此表}
    先按左列实现与测试，再按中列检查 QUDA 特有的状态更新；右列是不能从“V-cycle”三个字推出来、但会影响复现的行为。
  \end{block}
  \source{以上全部为当前快照的直接源码行号；伪代码并非源码直引，而是按控制流重写。}
\end{frame}

\begin{frame}[t]{源码控制流对照（2/2）：transfer、Galerkin、运行期和验证}
  \scriptsize
  \begin{tabularx}{\textwidth}{@{}l l X@{}}
    \toprule
    复现伪代码阶段 & QUDA 实现 & 对齐点 \\
    \midrule
    geometry map & \codepath{lib/transfer.cpp:200-241} & fine-to-coarse 逐点除 block；coarse-to-fine 按 coarse index 排序。 \\
    spin map & \codepath{lib/transfer.cpp:243-257} & non-staggered 为 \(s/b_s\)，staggered parity 映射到 0/1。 \\
    P/R dispatch & \codepath{lib/transfer.cpp:259-328} & aggregate、optimized KD、coarse KD 是不同 transfer；本报告主线是 aggregate。 \\
    block ortho & \codepath{include/kernels/block_orthogonalize.cuh:141-253} & aggregate 内 dot/reduce/subtract/normalize，支持重复和 tile。 \\
    UV/VUV & \codepath{lib/coarse_op.cuh:1043-1457} & nFace、双向、long link、diagonal split、identity/mass/twist/convert。 \\
    coarse apply & \codepath{include/kernels/dslash_coarse.cuh:126-324} & forward/backward local/halo gather，local \(X\)，\(-\kappa\) 汇总。 \\
    Yhat/Schur & \codepath{lib/dirac_coarse.cpp:506-649} & \(X^{-1}\)、\(\widehat Y\)、M、prepare、reconstruct、向下粗化 \(\widehat Y\)。 \\
    V-cycle & \codepath{lib/multigrid.cpp:1131-1224} & residual 来源和 solution type 分支。 \\
    verify & \codepath{lib/multigrid.cpp:745-1128} & 阈值、五类检查及 ASQTAD skip 条件。 \\
    \bottomrule
  \end{tabularx}
  \vspace{0.12cm}
  \begin{alertblock}{最终原则}
    先复现左列的数学动作，再逐项加入右列的布局/通信优化；不把某个 CUDA kernel 的线程映射误认为 multigrid 算法本身。
  \end{alertblock}
  \source{行号均来自 QUDA 1.1.0 参考源码快照；本页是索引，不替代对应源码阅读。}
\end{frame}

\appendix

% 附录源码按语义和长度手工拆页；每段最多 12 行，避免直引被裁切。
\begin{frame}[fragile,t]{附录 A1：运行期先确定 solution type（源码直引）}
  \lstinputlisting[firstline=1171,lastline=1180,basicstyle=\ttfamily\scriptsize]{../lib/multigrid.cpp}
  \begin{block}{对应伪代码}
    这是非底层分支的入口：创建工作字段，执行 \texttt{prepare}，再决定是否预平滑。
  \end{block}
  \source{原文件：\codepath{lib/multigrid.cpp:1171-1180}；代码未转写。}
\end{frame}

\begin{frame}[fragile,t]{附录 A2：显式 residual 分支（源码直引）}
  \lstinputlisting[firstline=1181,lastline=1190,basicstyle=\ttfamily\scriptsize]{../lib/multigrid.cpp}
  \begin{block}{对应伪代码}
    只有不一致的 smoother/coarse solution type 才显式重算 \(r=b-Ax\)；其余路径复用 solver residual 或输入。
  \end{block}
  \source{原文件：\codepath{lib/multigrid.cpp:1181-1190}；代码未转写。}
\end{frame}

\begin{frame}[fragile,t]{附录 A3：限制、粗解与回传（源码直引）}
  \lstinputlisting[firstline=1191,lastline=1200,basicstyle=\ttfamily\scriptsize]{../lib/multigrid.cpp}
  \begin{block}{对应伪代码}
    \texttt{R} 把本层 residual 送到下一层；粗解返回后按 full 或目标 parity 调用 \texttt{P}。
  \end{block}
  \source{原文件：\codepath{lib/multigrid.cpp:1191-1200}；代码未转写。}
\end{frame}

\begin{frame}[fragile,t]{附录 A4：有预平滑时必须做 correction add（源码直引）}
  \lstinputlisting[firstline=1201,lastline=1210,basicstyle=\ttfamily\scriptsize]{../lib/multigrid.cpp}
  \begin{alertblock}{不可替换语义}
    有 pre-smoother 时，\texttt{P} 的输出写到临时 correction，再执行 \(x\leftarrow x+Px_c\)；直接覆盖 \(x\) 会丢失平滑结果。
  \end{alertblock}
  \source{原文件：\codepath{lib/multigrid.cpp:1201-1210}；代码未转写。}
\end{frame}

\begin{frame}[fragile,t]{附录 A5：post smoother 与 coarsest-level 分支（源码直引）}
  \lstinputlisting[firstline=1211,lastline=1221,basicstyle=\ttfamily\scriptsize]{../lib/multigrid.cpp}
  \begin{block}{对应伪代码}
    回到本层后重新 \texttt{prepare}（若 solution type 改变），执行 post，再以 outer solution type \texttt{reconstruct}。
  \end{block}
  \source{原文件：\codepath{lib/multigrid.cpp:1211-1221}；代码未转写。}
\end{frame}

\begin{frame}[fragile,t]{附录 B1：块正交化的循环边界（源码直引）}
  \lstinputlisting[firstline=170,lastline=181,basicstyle=\ttfamily\scriptsize]{../include/kernels/block_orthogonalize.cuh}
  \begin{block}{对应伪代码}
    外层是重复次数，内层按 \(mVec\) 分批；第一次从 \(B\) 读，后续从已写回的 \(V\) 读。
  \end{block}
  \source{原文件：\codepath{include/kernels/block_orthogonalize.cuh:170-181}；代码未转写。}
\end{frame}

\begin{frame}[fragile,t]{附录 B2：按 chirality 读取候选列（源码直引）}
  \lstinputlisting[firstline=182,lastline=193,basicstyle=\ttfamily\scriptsize]{../include/kernels/block_orthogonalize.cuh}
  \begin{block}{对应伪代码}
    奇偶与 chirality 的索引在装载 \(B_i\) 时已经决定；这就是 spin map 不能在 P/R 阶段另起一套的原因。
  \end{block}
  \source{原文件：\codepath{include/kernels/block_orthogonalize.cuh:182-193}；代码未转写。}
\end{frame}

\begin{frame}[fragile,t]{附录 B3：跨历史列的块内积（源码直引）}
  \lstinputlisting[firstline=194,lastline=205,basicstyle=\ttfamily\scriptsize]{../include/kernels/block_orthogonalize.cuh}
  \begin{block}{对应伪代码}
    每个历史列 \(i<j\) 都先做 aggregate reduction；得到的复系数再用于整个当前 tile 的减法。
  \end{block}
  \source{原文件：\codepath{include/kernels/block_orthogonalize.cuh:194-205}；代码未转写。}
\end{frame}

\begin{frame}[fragile,t]{附录 B4：减去历史列并进入 tile 正交（源码直引）}
  \lstinputlisting[firstline=206,lastline=217,basicstyle=\ttfamily\scriptsize]{../include/kernels/block_orthogonalize.cuh}
  \begin{alertblock}{顺序}
    先减去 \(i<j\) 的历史列，再对 tile 内 \(m\) 列做列间正交；两层循环不能合并成一次局部归一。
  \end{alertblock}
  \source{原文件：\codepath{include/kernels/block_orthogonalize.cuh:206-217}；代码未转写。}
\end{frame}

\begin{frame}[fragile,t]{附录 B5：tile 内列间内积（源码直引）}
  \lstinputlisting[firstline=218,lastline=229,basicstyle=\ttfamily\scriptsize]{../include/kernels/block_orthogonalize.cuh}
  \begin{block}{对应伪代码}
    对当前列 \(m\)，只与 \(i<m\) 的当前 tile 列做内积，并把所有线程贡献归并到 aggregate 范围。
  \end{block}
  \source{原文件：\codepath{include/kernels/block_orthogonalize.cuh:218-229}；代码未转写。}
\end{frame}

\begin{frame}[fragile,t]{附录 B6：tile 内减法与归一化系数（源码直引）}
  \lstinputlisting[firstline=230,lastline=241,basicstyle=\ttfamily\scriptsize]{../include/kernels/block_orthogonalize.cuh}
  \begin{block}{对应伪代码}
    先按 \(i<m\) 更新 \(v_m\)，再用 aggregate norm 求 \(n_m^{-1}\)；零范数由实现保留为零系数并应在 reference 中报错。
  \end{block}
  \source{原文件：\codepath{include/kernels/block_orthogonalize.cuh:230-241}；代码未转写。}
\end{frame}

\begin{frame}[fragile,t]{附录 B7：归一化写回与重复循环结束（源码直引）}
  \lstinputlisting[firstline=242,lastline=253,basicstyle=\ttfamily\scriptsize]{../include/kernels/block_orthogonalize.cuh}
  \begin{alertblock}{实现边界}
    归一化后的 tile 才写入 \(V\)；下一批次和下一次重复都读取这个持久结果。
  \end{alertblock}
  \source{原文件：\codepath{include/kernels/block_orthogonalize.cuh:242-253}；代码未转写。}
\end{frame}

\begin{frame}[fragile,t]{附录 B8：VUV 的坐标与 aggregate 判定（源码直引）}
  \lstinputlisting[firstline=1608,lastline=1619,basicstyle=\ttfamily\scriptsize]{../include/kernels/coarse_op_kernel.cuh}
  \begin{block}{对应伪代码}
    先从 checkerboard 坐标恢复 coarse 坐标，再用 fine hop 的终点判断该矩阵块属于 \(X\) 还是 \(Y\)。
  \end{block}
  \source{原文件：\codepath{include/kernels/coarse_op_kernel.cuh:1608-1619}；代码未转写。}
\end{frame}

\begin{frame}[fragile,t]{附录 B9：VUV 的 diagonal 判定与 coarse parity（源码直引）}
  \lstinputlisting[firstline=1620,lastline=1631,basicstyle=\ttfamily\scriptsize]{../include/kernels/coarse_op_kernel.cuh}
  \begin{block}{对应伪代码}
    \texttt{isDiagonal} 决定 \(X/Y\)，coarse parity 与 checkerboard index 决定矩阵写回位置；这两个判断不能互换。
  \end{block}
  \source{原文件：\codepath{include/kernels/coarse_op_kernel.cuh:1620-1631}；代码未转写。}
\end{frame}

\begin{frame}[fragile,t]{附录 B10：VUV tile 乘法与局部系数（源码直引）}
  \lstinputlisting[firstline=1632,lastline=1643,basicstyle=\ttfamily\scriptsize]{../include/kernels/coarse_op_kernel.cuh}
  \begin{alertblock}{系数检查}
    \(VUV\) 先形成 coarse spin block；diagonal hopping 再乘 \(-\kappa\)，而 local \(X\) 的其他项另由专门 kernel 加入。
  \end{alertblock}
  \source{原文件：\codepath{include/kernels/coarse_op_kernel.cuh:1632-1643}；代码未转写。}
\end{frame}

\begin{frame}[fragile,t]{附录 B11：VUV 的 atomic 写回（源码直引）}
  \lstinputlisting[firstline=1644,lastline=1647,basicstyle=\ttfamily\scriptsize]{../include/kernels/coarse_op_kernel.cuh}
  \begin{block}{对应伪代码}
    最后依据共享或全局 atomic 路径写回；线程映射可以改变，\(X/Y\) 分类和矩阵转置语义不能改变。
  \end{block}
  \source{原文件：\codepath{include/kernels/coarse_op_kernel.cuh:1644-1647}；代码未转写。}
\end{frame}

\begin{frame}[fragile,t]{附录 C1：Yhat 先处理 backward link（源码直引）}
  \lstinputlisting[firstline=60,lastline=71,basicstyle=\ttfamily\scriptsize]{../include/kernels/coarse_op_preconditioned.cuh}
  \begin{block}{对应伪代码}
    边界使用 ghost index，内部使用 hop index；两条路径都加载同一个目标点的 \(X^{-\dagger}\)。
  \end{block}
  \source{原文件：\codepath{include/kernels/coarse_op_preconditioned.cuh:60-71}；代码未转写。}
\end{frame}

\begin{frame}[fragile,t]{附录 C2：Yhat backward 的矩阵乘与存储（源码直引）}
  \lstinputlisting[firstline=72,lastline=83,basicstyle=\ttfamily\scriptsize]{../include/kernels/coarse_op_preconditioned.cuh}
  \begin{alertblock}{方向}
    backward 的 \(Y\) 在邻点位置读取，并以 \(X^{-\dagger}\) 右乘；这正是 PC link 最容易写反的地方。
  \end{alertblock}
  \source{原文件：\codepath{include/kernels/coarse_op_preconditioned.cuh:72-83}；代码未转写。}
\end{frame}

\begin{frame}[fragile,t]{附录 C3：Yhat 内部 backward 路径（源码直引）}
  \lstinputlisting[firstline=84,lastline=95,basicstyle=\ttfamily\scriptsize]{../include/kernels/coarse_op_preconditioned.cuh}
  \begin{block}{对应伪代码}
    非边界使用 \texttt{back\_idx}；矩阵 tile 的 \(B^T\) 装载实现右乘，不能替换为普通左乘。
  \end{block}
  \source{原文件：\codepath{include/kernels/coarse_op_preconditioned.cuh:84-95}；代码未转写。}
\end{frame}

\begin{frame}[fragile,t]{附录 C4：Yhat backward 的输出（源码直引）}
  \lstinputlisting[firstline=96,lastline=107,basicstyle=\ttfamily\scriptsize]{../include/kernels/coarse_op_preconditioned.cuh}
  \begin{block}{对应伪代码}
    backward 结果保存到反向方向的 link；随后进入 forward 的 \(X^{-1}Y\) 路径。
  \end{block}
  \source{原文件：\codepath{include/kernels/coarse_op_preconditioned.cuh:96-107}；代码未转写。}
\end{frame}

\begin{frame}[fragile,t]{附录 C5：Yhat forward 是左乘（源码直引）}
  \lstinputlisting[firstline=108,lastline=119,basicstyle=\ttfamily\scriptsize]{../include/kernels/coarse_op_preconditioned.cuh}
  \begin{alertblock}{方向}
    forward 先加载 \(X^{-1}(q)\)，再加载 \(Y_{+\mu}(q)\)，计算 \(X^{-1}Y\) 并保存到正向方向。
  \end{alertblock}
  \source{原文件：\codepath{include/kernels/coarse_op_preconditioned.cuh:108-119}；代码未转写。}
\end{frame}

\begin{frame}[fragile,t]{附录 C6：Yhat forward 写回（源码直引）}
  \lstinputlisting[firstline=120,lastline=121,basicstyle=\ttfamily\scriptsize]{../include/kernels/coarse_op_preconditioned.cuh}
  \begin{block}{对应伪代码}
    两个方向的 ghost/内部索引都在本 kernel 内闭合；外部实现只需保持相同的矩阵方向和字段布局。
  \end{block}
  \source{原文件：\codepath{include/kernels/coarse_op_preconditioned.cuh:120-121}；代码未转写。}
\end{frame}

\begin{frame}[fragile,t]{附录 D1：coarse dslash 的 forward gather（源码直引）}
  \lstinputlisting[firstline=148,lastline=159,basicstyle=\ttfamily\scriptsize]{../include/kernels/dslash_coarse.cuh}
  \begin{block}{对应伪代码}
    forward 访问 \(q+\hat\mu\)；跨 rank 时把本地 index 换成 ghost index，目标 parity 保持不变。
  \end{block}
  \source{原文件：\codepath{include/kernels/dslash_coarse.cuh:148-159}；代码未转写。}
\end{frame}

\begin{frame}[fragile,t]{附录 D2：forward halo 中的矩阵乘（源码直引）}
  \lstinputlisting[firstline=160,lastline=171,basicstyle=\ttfamily\scriptsize]{../include/kernels/dslash_coarse.cuh}
  \begin{block}{对应伪代码}
    ghost spinor 按 source batch、spin、color 索引；\(Y_{+\mu}\) 作用在输出 spin/color 行上。
  \end{block}
  \source{原文件：\codepath{include/kernels/dslash_coarse.cuh:160-171}；代码未转写。}
\end{frame}

\begin{frame}[fragile,t]{附录 D3：forward bulk 与 dagger 分支（源码直引）}
  \lstinputlisting[firstline=172,lastline=183,basicstyle=\ttfamily\scriptsize]{../include/kernels/dslash_coarse.cuh}
  \begin{alertblock}{核对点}
    dagger 模式改变读取的方向 link；非 dagger 模式读取 \(Y[d+4]\)。这不是简单地给最终向量取共轭。
  \end{alertblock}
  \source{原文件：\codepath{include/kernels/dslash_coarse.cuh:172-183}；代码未转写。}
\end{frame}

\begin{frame}[fragile,t]{附录 D4：forward bulk 读取（源码直引）}
  \lstinputlisting[firstline=184,lastline=195,basicstyle=\ttfamily\scriptsize]{../include/kernels/dslash_coarse.cuh}
  \begin{block}{对应伪代码}
    内部点直接用 \texttt{fwd\_idx} 读取输入；通信和计算的数学结果必须逐点一致。
  \end{block}
  \source{原文件：\codepath{include/kernels/dslash_coarse.cuh:184-195}；代码未转写。}
\end{frame}

\begin{frame}[fragile,t]{附录 D5：切换到 backward gather（源码直引）}
  \lstinputlisting[firstline=196,lastline=207,basicstyle=\ttfamily\scriptsize]{../include/kernels/dslash_coarse.cuh}
  \begin{block}{对应伪代码}
    backward 使用 \(q-\hat\mu\) 的 link 和输入；GPU 中 forward/backward 可由不同线程负责，但归约顺序需固定。
  \end{block}
  \source{原文件：\codepath{include/kernels/dslash_coarse.cuh:196-207}；代码未转写。}
\end{frame}

\begin{frame}[fragile,t]{附录 D6：backward 的 ghost 分支（源码直引）}
  \lstinputlisting[firstline=208,lastline=219,basicstyle=\ttfamily\scriptsize]{../include/kernels/dslash_coarse.cuh}
  \begin{alertblock}{边界}
    backward 边界先用 \texttt{ghostFaceIndex} 定位 link 与 spinor；两者的 face 顺序必须使用同一 parity 约定。
  \end{alertblock}
  \source{原文件：\codepath{include/kernels/dslash_coarse.cuh:208-219}；代码未转写。}
\end{frame}

\begin{frame}[fragile,t]{附录 D7：backward link 的共轭转置（源码直引）}
  \lstinputlisting[firstline=220,lastline=231,basicstyle=\ttfamily\scriptsize]{../include/kernels/dslash_coarse.cuh}
  \begin{block}{对应伪代码}
    非 dagger 路径使用 \(Y_d^\dagger(q-\hat\mu)\)；这一共轭转置属于 link，不属于输入 spinor。
  \end{block}
  \source{原文件：\codepath{include/kernels/dslash_coarse.cuh:220-231}；代码未转写。}
\end{frame}

\begin{frame}[fragile,t]{附录 D8：backward bulk 与局部 link 索引（源码直引）}
  \lstinputlisting[firstline=232,lastline=243,basicstyle=\ttfamily\scriptsize]{../include/kernels/dslash_coarse.cuh}
  \begin{block}{对应伪代码}
    bulk 路径使用 \texttt{back\_idx}；\(Y\) 的 parity 是邻点 parity，而输出仍属于当前目标 parity。
  \end{block}
  \source{原文件：\codepath{include/kernels/dslash_coarse.cuh:232-243}；代码未转写。}
\end{frame}

\begin{frame}[fragile,t]{附录 D9：local X 与 \(-\kappa\) 汇总（源码直引）}
  \lstinputlisting[firstline=244,lastline=249,basicstyle=\ttfamily\scriptsize]{../include/kernels/dslash_coarse.cuh}
  \begin{block}{对应伪代码}
    hopping 汇总后才统一乘 \(-\kappa\)；\(X\) 作为 local matrix 另行加入，故装配时不能把该系数重复乘进 \(Y\)。
  \end{block}
  \source{原文件：\codepath{include/kernels/dslash_coarse.cuh:244-249}；代码未转写。}
\end{frame}

\begin{frame}[fragile,t]{附录 D：建议的最小配置与报告记录模板}
  \begin{columns}[T,totalwidth=\textwidth]
    \begin{column}{0.47\textwidth}
      \begin{lstlisting}[firstnumber=1]
geometry: Lf[4], periodic/BC
operator: Wilson or Clover, kappa, mass
level[0]: b[4], spin_block, nvec
setup: num_setup_iter, maxiter, tol
ortho: global_pre/post, n_block_ortho
solve: smoother, nu_pre, nu_post, omega
coarse: solution_type, solve_type, cycle
precision: B/V/Y/X/halo
parallel: nranks, rank_grid, face depth
verify: V1..V5 values and thresholds
      \end{lstlisting}
    \end{column}
    \begin{column}{0.49\textwidth}
      \begin{block}{每次回归至少记录}
        \begin{tabularx}{\linewidth}{@{}lX@{}}
          \toprule
          字段 & 内容 \\
          \midrule
          输入 & gauge/clover 来源、随机种子、边界 \\
          结构 & 每层 (L,b,Nvec,Nspin,parity) \\
          算法 & setup type、smoother、cycle、PC \\
          误差 & V1--V5 的 L2/max、真 residual \\
          性能 & 设备、rank、时间、基线（若已测） \\
          \bottomrule
        \end{tabularx}
      \end{block}
      \begin{alertblock}{真实性}
        没有动态命令和输出时，性能/收敛栏写“未验证”，不能用源码行号代替运行证据。
      \end{alertblock}
    \end{column}
  \end{columns}
  \source{参数字段：\codepath{include/quda.h:638-835}；测试参数入口：\codepath{tests/utils/set_params.cpp:361-627}。本页是复现实验记录模板。}
\end{frame}

\begin{frame}[t]{最终结论：复现 QUDA multigrid 的最短正确路径}
  \begin{center}
    \begin{tikzpicture}[node distance=0.22cm and 0.20cm]
      \node[flowwide] (a) {固定坐标、parity、field layout};
      \node[flowwide,right=of a] (b) {生成并块正交化 \(B\to V\)};
      \node[flowwide,right=of b] (c) {实现 \(P/R\)，通过 \(RP\)};
      \node[flowwide,right=of c] (d) {用显式 \(RDP\) 得 \(D_c\)};
      \node[flowwide,below=0.36cm of d] (e) {拆 (Y/X)，通过 Galerkin};
      \node[flowwide,left=of e] (f) {加 \(X^{-1}/\widehat Y\)，通过 PC};
      \node[flowwide,left=of f] (g) {加 halo、MR、recursive V-cycle};
      \draw[arrow] (a)--(b)--(c)--(d);
      \draw[arrow] (d)--(e)--(f)--(g);
    \end{tikzpicture}
  \end{center}
  \vspace{0.16cm}
  \begin{columns}[T,totalwidth=\textwidth]
    \begin{column}{0.31\textwidth}
      \begin{block}{已完成}
        本版补齐语言无关 setup、null-vector batch、两级块正交、P/R、Galerkin、PC Schur、halo、MR、GCR、参数映射和复现阶梯。
      \end{block}
    \end{column}
    \begin{column}{0.31\textwidth}
      \begin{block}{可证据化}
        每个关键步骤均有 QUDA 文件:行号；附录保留运行期、正交化、VUV、Yhat、dslash 直引片段。
      \end{block}
    \end{column}
    \begin{column}{0.31\textwidth}
      \begin{alertblock}{尚待动态验证}
        单 GPU/MPI 实际残差、不同参数的收敛曲线和硬件性能；应按“单 rank reference→GPU→MPI→多层”阶梯补测。
      \end{alertblock}
    \end{column}
  \end{columns}
  \source{总结依据：\codepath{lib/multigrid.cpp:1131-1224}、\codepath{lib/transfer.cpp:200-328}、\codepath{lib/dirac_coarse.cpp:121-170, 506-649}；结论中的动态缺口为本次工作边界。}
\end{frame}

\end{document}
